\input{preamble}

% OK, start here.
%
\begin{document}

\title{Descente et espaces algébriques}

\maketitle

\phantomsection
\label{section-phantom}

\tableofcontents

\section{Introduction}
\label{section-introduction}

\noindent
Dans le chapitre consacré aux topologies sur les espaces algébriques (voir
Topologies sur les espaces, section \ref{spaces-topologies-section-introduction})
nous avons introduit les recouvrements \'etales, fppf, lisses, syntomiques et fpqc des
espaces algébriques.
Dans ce chapitre, nous examinons quelles structures sur les espaces algébriques
peuvent être descendues le long de tels recouvrements.
Voir par exemple \cite{Gr-I}, \cite{Gr-II}, \cite{Gr-III},
\cite{Gr-IV}, \cite{Gr-V} et \cite{Gr-VI}.



\section{Conventions}
\label{section-conventions}

\noindent
L'hypothèse générale est que tous les schémas sont contenus dans
un grand site fppf $\Sch_{fppf}$. Tous les anneaux $A$ considérés
ont la propriété que $\Spec(A)$ est (isomorphe à) un
objet de ce grand site.

\medskip\noindent
Soit $S$ un schéma et soit $X$ un espace algébrique sur $S$.
Dans ce chapitre et dans le suivant, nous écrirons $X \times_S X$
pour le produit de $X$ par lui-même (dans la catégorie des espaces algébriques
sur $S$), au lieu de $X \times X$.






\section{Données de descente pour les faisceaux quasi-cohérents}
\label{section-equivalence}

\noindent
Cette section est l'analogue de
Descente, section \ref{descent-section-equivalence}
pour les espaces algébriques.
Il est utile de lire d'abord cette section.

\begin{definition}
\label{definition-descent-datum-quasi-coherent}
Soit $S$ un schéma. Soit $\{f_i : X_i \to X\}_{i \in I}$ une famille
de morphismes d'espaces algébriques sur $S$ de but fixé $X$.
\begin{enumerate}
\item Une {\it donnée de descente $(\mathcal{F}_i, \varphi_{ij})$
pour les faisceaux quasi-cohérents} relativement à la famille donnée
est constituée d'un faisceau quasi-cohérent $\mathcal{F}_i$ sur $X_i$ pour
chaque $i \in I$, d'un isomorphisme de
$\mathcal{O}_{X_i \times_X X_j}$-modules quasi-cohérents
$\varphi_{ij} : \text{pr}_0^*\mathcal{F}_i \to \text{pr}_1^*\mathcal{F}_j$
pour chaque paire $(i, j) \in I^2$
tel que, pour tout triplet d'indices $(i, j, k) \in I^3$, le
diagramme
$$
\xymatrix{
\text{pr}_0^*\mathcal{F}_i \ar[rd]_{\text{pr}_{01}^*\varphi_{ij}}
\ar[rr]_{\text{pr}_{02}^*\varphi_{ik}} & &
\text{pr}_2^*\mathcal{F}_k \\
& \text{pr}_1^*\mathcal{F}_j \ar[ru]_{\text{pr}_{12}^*\varphi_{jk}} &
}
$$
de $\mathcal{O}_{X_i \times_X X_j \times_X X_k}$-modules
soit commutatif. C'est la {\it condition de cocycle}.
\item Un {\it morphisme $\psi : (\mathcal{F}_i, \varphi_{ij}) \to
(\mathcal{F}'_i, \varphi'_{ij})$ de données de descente} est donné
par une famille $\psi = (\psi_i)_{i\in I}$ de morphismes de
$\mathcal{O}_{X_i}$-modules $\psi_i : \mathcal{F}_i \to \mathcal{F}'_i$
tels que tous les diagrammes
$$
\xymatrix{
\text{pr}_0^*\mathcal{F}_i \ar[r]_{\varphi_{ij}} \ar[d]_{\text{pr}_0^*\psi_i}
& \text{pr}_1^*\mathcal{F}_j \ar[d]^{\text{pr}_1^*\psi_j} \\
\text{pr}_0^*\mathcal{F}'_i \ar[r]^{\varphi'_{ij}} &
\text{pr}_1^*\mathcal{F}'_j \\
}
$$
soient commutatifs.
\end{enumerate}
\end{definition}

\begin{lemma}
\label{lemma-map-families}
Soit $S$ un schéma.
Soient $\mathcal{U} = \{U_i \to U\}_{i \in I}$ et
$\mathcal{V} = \{V_j \to V\}_{j \in J}$
des familles de morphismes d'espaces algébriques sur $S$ à buts fixés.
Soit $(g, \alpha : I \to J, (g_i)) : \mathcal{U} \to \mathcal{V}$
un morphisme de familles de morphismes à but fixé, voir
Sites, définition \ref{sites-definition-morphism-coverings}.
Soit $(\mathcal{F}_j, \varphi_{jj'})$ une donnée de descente
pour les faisceaux quasi-cohérents relativement à la
famille $\{V_j \to V\}_{j \in J}$. Alors
\begin{enumerate}
\item Le système
$$
\left(g_i^*\mathcal{F}_{\alpha(i)},
(g_i \times g_{i'})^*\varphi_{\alpha(i)\alpha(i')}\right)
$$
est une donnée de descente relativement à la famille $\{U_i \to U\}_{i \in I}$.
\item Cette construction est fonctorielle en la donnée de descente
$(\mathcal{F}_j, \varphi_{jj'})$.
\item Étant donné un second morphisme
$(g', \alpha' : I \to J, (g'_i))$
de familles de morphismes à but fixé
tel que $g = g'$, il existe un isomorphisme fonctoriel de données de descente
$$
(g_i^*\mathcal{F}_{\alpha(i)},
(g_i \times g_{i'})^*\varphi_{\alpha(i)\alpha(i')})
\cong
((g'_i)^*\mathcal{F}_{\alpha'(i)},
(g'_i \times g'_{i'})^*\varphi_{\alpha'(i)\alpha'(i')}).
$$
\end{enumerate}
\end{lemma}

\begin{proof}
Démonstration omise. Indication : les morphismes
$g_i^*\mathcal{F}_{\alpha(i)} \to (g'_i)^*\mathcal{F}_{\alpha'(i)}$
qui donnent l'isomorphisme des données de descente en (3)
sont les images réciproques des morphismes $\varphi_{\alpha(i)\alpha'(i)}$ par les
morphismes $(g_i, g'_i) : U_i \to V_{\alpha(i)} \times_V V_{\alpha'(i)}$.
\end{proof}

\noindent
Soit $g : U \to V$ un morphisme d'espaces algébriques.
Le lemme ci-dessus nous indique qu'il existe un foncteur d'image réciproque bien défini
entre les catégories de données de descente relatives à des familles de
morphismes de but $V$ et $U$, pourvu qu'il existe un morphisme entre ces
familles de morphismes qui ``vit au-dessus de $g$''.

\begin{definition}
\label{definition-descent-datum-effective-quasi-coherent}
Soit $S$ un schéma.
Soit $\{U_i \to U\}_{i \in I}$ une famille de morphismes d'espaces algébriques
sur $S$ de but fixé.
\begin{enumerate}
\item Soit $\mathcal{F}$ un $\mathcal{O}_U$-module quasi-cohérent.
Nous appelons l'unique donnée de descente sur $\mathcal{F}$ relativement au recouvrement
$\{U \to U\}$ la {\it donnée de descente triviale}.
\item L'image réciproque de la donnée de descente triviale sur
$\{U_i \to U\}$ est appelée la {\it donnée de descente canonique}.
Notation : $(\mathcal{F}|_{U_i}, can)$.
\item Une donnée de descente $(\mathcal{F}_i, \varphi_{ij})$
pour les faisceaux quasi-cohérents relativement à la famille donnée
est dite {\it effective} s'il existe un faisceau quasi-cohérent
$\mathcal{F}$ sur $U$ tel que $(\mathcal{F}_i, \varphi_{ij})$
soit isomorphe à $(\mathcal{F}|_{U_i}, can)$.
\end{enumerate}
\end{definition}

\begin{lemma}
\label{lemma-zariski-descent-effective}
Soit $S$ un schéma. Soit $U$ un espace algébrique sur $S$.
Soit $\{U_i \to U\}$ un recouvrement de Zariski de $U$, voir
Topologies sur les espaces,
définition \ref{spaces-topologies-definition-zariski-covering}.
Toute donnée de descente sur des faisceaux quasi-cohérents
relativement à la famille $\mathcal{U} = \{U_i \to U\}$ est
effective. De plus, le foncteur de la catégorie des
$\mathcal{O}_U$-modules quasi-cohérents vers la catégorie
des données de descente relativement à $\{U_i \to U\}$ est pleinement fidèle.
\end{lemma}

\begin{proof}
Démonstration omise.
\end{proof}












\section{Descente fpqc des faisceaux quasi-cohérents}
\label{section-fpqc-descent-quasi-coherent}

\noindent
La principale application de la descente plate pour les modules est
l'énoncé de descente correspondant pour les faisceaux quasi-cohérents
relativement aux recouvrements fpqc.

\begin{proposition}
\label{proposition-fpqc-descent-quasi-coherent}
Soit $S$ un schéma.
Soit $\{X_i \to X\}$ un recouvrement fpqc d'espaces algébriques sur $S$, voir
Topologies sur les espaces,
définition \ref{spaces-topologies-definition-fpqc-covering}.
Toute donnée de descente sur des faisceaux quasi-cohérents
relativement à $\{X_i \to X\}$ est effective.
De plus, le foncteur de la catégorie des
$\mathcal{O}_X$-modules quasi-cohérents vers la catégorie
des données de descente relativement à $\{X_i \to X\}$ est pleinement fidèle.
\end{proposition}

\begin{proof}
Ceci est plus ou moins une conséquence formelle du
résultat correspondant pour les schémas, voir
Descente, proposition \ref{descent-proposition-fpqc-descent-quasi-coherent}.
Voici une stratégie de démonstration :
\begin{enumerate}
\item Le fait que $\{X_i \to X\}$ soit un raffinement du recouvrement trivial
$\{X \to X\}$ donne, via
le lemme \ref{lemma-map-families},
un foncteur $\QCoh(\mathcal{O}_X) \to DD(\{X_i \to X\})$ de la
catégorie des $\mathcal{O}_X$-modules quasi-cohérents vers la catégorie des
données de descente pour la famille donnée.
\item Afin de démontrer la proposition, nous allons construire un
foncteur quasi-inverse
$back : DD(\{X_i \to X\}) \to \QCoh(\mathcal{O}_X)$.
\item En appliquant à nouveau
le lemme \ref{lemma-map-families}
nous voyons qu'il existe un foncteur
$DD(\{X_i \to X\}) \to DD(\{T_j \to X\})$
si $\{T_j \to X\}$ est un raffinement de la famille donnée.
Par conséquent, pour construire le foncteur $back$, nous pouvons supposer que
chaque $X_i$ est un schéma, voir
Topologies sur les espaces,
lemme \ref{spaces-topologies-lemma-refine-fpqc-schemes}.
Cela nous ramène au cas où tous les $X_i$ sont des schémas.
\item Un faisceau quasi-cohérent sur $X$ est, par définition, un
$\mathcal{O}_X$-module quasi-cohérent sur $X_\etale$. Pour tout
$U \in \Ob(X_\etale)$, on obtient alors un recouvrement fppf
$\{U_i \times_X X_i \to U\}$ par des schémas et un morphisme
$g : \{U_i \times_X X_i \to U\} \to \{X_i \to X\}$ de recouvrements
au-dessus de $U \to X$. Étant donnée une donnée de descente
$\xi = (\mathcal{F}_i, \varphi_{ij})$, on obtient un
$\mathcal{O}_U$-module quasi-cohérent $\mathcal{F}_{\xi, U}$ correspondant
à l'image réciproque $g^*\xi$ fournie par le
lemme \ref{lemma-map-families}
sur le recouvrement de $U$, en utilisant l'effectivité des recouvrements fppf de schémas, voir
Descente, proposition \ref{descent-proposition-fpqc-descent-quasi-coherent}.
\item Vérifier que $\xi \mapsto \mathcal{F}_{\xi, U}$ est fonctorielle en $\xi$.
Démonstration omise.
\item Vérifier que $\xi \mapsto \mathcal{F}_{\xi, U}$ est compatible
avec les morphismes $U \to U'$ du site $X_\etale$, de sorte que
le système de faisceaux $\mathcal{F}_{\xi, U}$ corresponde à un faisceau quasi-cohérent
$\mathcal{F}_\xi$ sur $X_\etale$, voir
Propriétés des espaces,
lemme \ref{spaces-properties-lemma-characterize-quasi-coherent-small-etale}.
Détails omis.
\item Vérifier que $back : \xi \mapsto \mathcal{F}_\xi$ est quasi-inverse
du foncteur construit en (1). Démonstration omise.
\end{enumerate}
Ceci achève la démonstration.
\end{proof}









\section{Modules quasi-cohérents et affines}
\label{section-alternative-quasi-coherent}

\noindent
Soit $S$ un schéma. Soit $X$ un espace algébrique sur $S$.
Rappelons que $X_{affine, \etale}$ est la sous-catégorie pleine de $X_\etale$
dont les objets sont affines, que l'on munit d'une structure de site en déclarant les recouvrements
comme étant les recouvrements \'etale standards. Voir Propriétés des espaces, définition
\ref{spaces-properties-definition-affine-etale-site}.
D'après Propriétés des espaces, lemme \ref{spaces-properties-lemma-alternative}
nous avons une équivalence de topos
$g : \Sh(X_{affine, \etale}) \to \Sh(X_\etale)$
dont le foncteur d'image réciproque est donné par la restriction.
Rappelons que $\mathcal{O}_X$ désigne le faisceau structural sur
$X_\etale$. Nous obtenons alors une équivalence
\begin{equation}
\label{equation-alternative-small-ringed}
(\Sh(X_{affine, \etale}), \mathcal{O}_X|_{X_{affine, \etale}})
\longrightarrow
(\Sh(X_\etale), \mathcal{O}_X)
\end{equation}
de topos annelés. Nous écrirons souvent $\mathcal{O}_X$
au lieu de $\mathcal{O}_X|_{X_{affine, \etale}}$.
Cela étant dit, nous pouvons également comparer les modules quasi-cohérents.

\begin{lemma}
\label{lemma-quasi-coherent-alternative-small}
Soit $S$ un schéma. Soit $X$ un espace algébrique sur $S$.
Soit $\mathcal{F}$ un préfaisceau de $\mathcal{O}_X$-modules
sur $X_{affine, \etale}$. Les conditions suivantes sont équivalentes
\begin{enumerate}
\item pour tout morphisme $U \to U'$ de $X_{affine, \etale}$ la flèche
$\mathcal{F}(U') \otimes_{\mathcal{O}_X(U')} \mathcal{O}_X(U)
\to \mathcal{F}(U)$ est un isomorphisme,
\item $\mathcal{F}$ est un module quasi-cohérent sur le site annelé
$(X_{affine, \etale}, \mathcal{O}_X)$ au sens de
Modules sur les sites, définition \ref{sites-modules-definition-site-local},
\item $\mathcal{F}$ correspond à un module quasi-cohérent sur $X$
via l'équivalence (\ref{equation-alternative-small-ringed}),
\end{enumerate}
\end{lemma}

\begin{proof}
Supposons (1) satisfaite. Pour montrer que $\mathcal{F}$ est un faisceau, soit
$\mathcal{U} = \{U_i \to U\}_{i = 1, \ldots, n}$ un recouvrement
de $X_{affine, \etale}$. La condition de faisceau pour $\mathcal{F}$
et $\mathcal{U}$, par notre hypothèse sur $\mathcal{F}$,
se réduit à montrer que
$$
0 \to \mathcal{F}(U) \to
\prod \mathcal{F}(U) \otimes_{\mathcal{O}_X(U)} \mathcal{O}_X(U_i) \to
\prod \mathcal{F}(U) \otimes_{\mathcal{O}_X(U)} \mathcal{O}_X(U_i \times_U U_j)
$$
est exacte. Cela résulte de ce que $\mathcal{O}_X(U) \to \prod \mathcal{O}_X(U_i)$
est fidèlement plat
(par Descente, lemme \ref{descent-lemma-standard-covering-Cech} et
le fait que les recouvrements dans $X_{affine, \etale}$ sont des recouvrements \'etale
standards) et nous pouvons appliquer Descente, lemme \ref{descent-lemma-ff-exact}.
Ensuite, nous montrons que $\mathcal{F}$ est quasi-cohérent sur $X_{affine, \etale}$.
En effet, pour $U$ dans $X_{affine, \etale}$, posons $R = \mathcal{O}_X(U)$
et choisissons une présentation
$$
\bigoplus\nolimits_{k \in K} R
\longrightarrow
\bigoplus\nolimits_{l \in L} R
\longrightarrow
\mathcal{F}(U)
\longrightarrow 0
$$
par des $R$-modules libres. D'après la propriété (1) et l'exactitude à droite du produit tensoriel
nous voyons que pour tout morphisme $U' \to U$ dans $X_{affine, \etale}$
nous obtenons une présentation
$$
\bigoplus\nolimits_{k \in K} \mathcal{O}_X(U')
\longrightarrow
\bigoplus\nolimits_{l \in L} \mathcal{O}_X(U')
\longrightarrow
\mathcal{F}(U')
\longrightarrow 0
$$
En d'autres termes, la restriction de $\mathcal{F}$
à la catégorie localisée $X_{affine, etale}/U$ possède une présentation
$$
\bigoplus\nolimits_{k \in K} \mathcal{O}_X|_{X_{affine, \etale}/U}
\longrightarrow
\bigoplus\nolimits_{l \in L} \mathcal{O}_X|_{X_{affine, \etale}/U}
\longrightarrow
\mathcal{F}|_{X_{affine, \etale}/U}
\longrightarrow 0
$$
comme requis pour montrer que $\mathcal{F}$ est quasi-cohérent.
Malgré cette notation quelque peu horrible, ceci achève la démonstration
que (1) implique (2).

\medskip\noindent
Puisque la notion de module quasi-cohérent est intrinsèque
(Modules sur les sites, lemme \ref{sites-modules-lemma-special-locally-free})
nous voyons que l'équivalence (\ref{equation-alternative-small-ringed})
induit une équivalence entre les catégories de modules quasi-cohérents.
Nous obtenons ainsi l'équivalence de (2) et (3).

\medskip\noindent
Supposons (3) et démontrons (1). Soit
$\mathcal{G}$ soit un module quasi-cohérent sur $X$ correspondant à $\mathcal{F}$.
Soit $h : U \to U' \to X$ un morphisme de $X_{affine, \etale}$.
Notons $f : U \to X$ et $f' : U' \to X$ les morphismes structuraux,
de sorte que $f = f' \circ h$. Nous avons
$\mathcal{F}(U') = \Gamma(U', (f')^*\mathcal{G})$ et
$\mathcal{F}(U) = \Gamma(U, f^*\mathcal{G}) = \Gamma(U, h^*(f')^*\mathcal{G})$.
La propriété (1) résulte alors de Schémas, lemme \ref{schemes-lemma-widetilde-pullback}.
\end{proof}





\section{Descente des propriétés de finitude des modules}
\label{section-descent-finiteness}

\noindent
Cette section est l'analogue, pour les espaces algébriques, de
Descente, section \ref{descent-section-descent-finiteness}.
Le but est de montrer qu'on peut vérifier qu'un module quasi-cohérent
satisfait certaines conditions de finitude en les vérifiant sur les membres d'un
recouvrement. Nous utiliserons à plusieurs reprises le schéma de preuve suivant.
Supposons que $X$ soit un espace algébrique et que $\{X_i \to X\}$
soit un recouvrement fppf (resp.\ fpqc). Soit $U \to X$ un morphisme \'etale
surjectif tel que $U$ soit un schéma. Alors il existe un
recouvrement fppf (resp.\ fpqc) $\{Y_j \to X\}$ tel que
\begin{enumerate}
\item $\{Y_j \to X\}$ est un raffinement de $\{X_i \to X\}$,
\item chaque $Y_j$ est un schéma, et
\item chaque morphisme $Y_j \to X$ se factorise par $U$, et
\item $\{Y_j \to U\}$ est un recouvrement fppf (resp.\ fpqc) de $U$.
\end{enumerate}
En effet, raffinons d'abord $\{X_i \to X\}$ par un recouvrement fppf (resp.\ fpqc)
tel que chaque $X_i$ soit un schéma, voir
Topologies sur les espaces, lemme \ref{spaces-topologies-lemma-refine-fppf-schemes},
resp.\ le lemme \ref{spaces-topologies-lemma-refine-fpqc-schemes}.
Posons alors $Y_i = U \times_X X_i$. Un module quasi-cohérent
sur $\mathcal{O}_X$, $\mathcal{F}$ est de type fini, de
présentation finie, etc. si et seulement si le module quasi-cohérent
sur $\mathcal{O}_U$, $\mathcal{F}|_U$, est de type fini, de
présentation finie, etc. Ainsi, l'existence du
raffinement $\{Y_j \to X\}$ permet de ramener la démonstration des
lemmes suivants au cas des schémas. Nous l'indiquerons en disant
que « {\it le résultat découle du cas des schémas par
localisation \'etale} ».

\begin{lemma}
\label{lemma-finite-type-descends}
Soit $X$ un espace algébrique au-dessus d'un schéma $S$.
Soit $\mathcal{F}$ un module quasi-cohérent sur $\mathcal{O}_X$.
Soit $\{f_i : X_i \to X\}_{i \in I}$ un recouvrement fpqc tel que
chaque $f_i^*\mathcal{F}$ soit un module de type fini sur $\mathcal{O}_{X_i}$.
Alors $\mathcal{F}$ est un module de type fini sur $\mathcal{O}_X$.
\end{lemma}

\begin{proof}
Cela découle du cas des schémas, voir
Descente, lemme \ref{descent-lemma-finite-type-descends},
par localisation \'etale.
\end{proof}

\begin{lemma}
\label{lemma-finite-presentation-descends}
Soit $X$ un espace algébrique au-dessus d'un schéma $S$.
Soit $\mathcal{F}$ un module quasi-cohérent sur $\mathcal{O}_X$.
Soit $\{f_i : X_i \to X\}_{i \in I}$ un recouvrement fpqc tel que
chaque $f_i^*\mathcal{F}$ soit un module sur $\mathcal{O}_{X_i}$ de présentation
finie. Alors $\mathcal{F}$ est un module sur $\mathcal{O}_X$
de présentation finie.
\end{lemma}

\begin{proof}
Cela découle du cas des schémas, voir
Descente, lemme \ref{descent-lemma-finite-presentation-descends},
par localisation \'etale.
\end{proof}

\begin{lemma}
\label{lemma-flat-descends}
Soit $X$ un espace algébrique au-dessus d'un schéma $S$.
Soit $\mathcal{F}$ un module quasi-cohérent sur $\mathcal{O}_X$.
Soit $\{f_i : X_i \to X\}_{i \in I}$ un recouvrement fpqc tel que
chaque $f_i^*\mathcal{F}$ soit un module plat sur $\mathcal{O}_{X_i}$.
Alors $\mathcal{F}$ est un module plat sur $\mathcal{O}_X$.
\end{lemma}

\begin{proof}
Cela découle du cas des schémas, voir
Descente, lemme \ref{descent-lemma-flat-descends},
par localisation \'etale.
\end{proof}

\begin{lemma}
\label{lemma-finite-locally-free-descends}
Soit $X$ un espace algébrique au-dessus d'un schéma $S$.
Soit $\mathcal{F}$ un module quasi-cohérent sur $\mathcal{O}_X$.
Soit $\{f_i : X_i \to X\}_{i \in I}$ un recouvrement fpqc tel que
chaque $f_i^*\mathcal{F}$ soit un module localement libre de rang fini sur $\mathcal{O}_{X_i}$.
Alors $\mathcal{F}$ est un module localement libre de rang fini sur $\mathcal{O}_X$.
\end{lemma}

\begin{proof}
Cela découle du cas des schémas, voir
Descente, lemme \ref{descent-lemma-finite-locally-free-descends},
par localisation \'etale.
\end{proof}

\noindent
La définition d'un faisceau quasi-cohérent localement projectif se trouve dans
Propriétés des espaces, section
\ref{spaces-properties-section-locally-projective}.
Il y est également démontré que cette notion est préservée par image réciproque.

\begin{lemma}
\label{lemma-locally-projective-descends}
Soit $X$ un espace algébrique au-dessus d'un schéma $S$.
Soit $\mathcal{F}$ un module quasi-cohérent sur $\mathcal{O}_X$.
Soit $\{f_i : X_i \to X\}_{i \in I}$ un recouvrement fpqc tel que
chaque $f_i^*\mathcal{F}$ soit un module localement projectif sur $\mathcal{O}_{X_i}$.
Alors $\mathcal{F}$ est un module localement projectif sur $\mathcal{O}_X$.
\end{lemma}

\begin{proof}
Cela découle du cas des schémas, voir
Descente, lemme \ref{descent-lemma-locally-projective-descends},
par localisation \'etale.
\end{proof}

\noindent
Nous ajoutons ici deux résultats liés aux précédents, mais
de nature légèrement différente.

\begin{lemma}
\label{lemma-finite-over-finite-module}
Soit $S$ un schéma.
Soit $f : X \to Y$ un morphisme d'espaces algébriques sur $S$.
Soit $\mathcal{F}$ un module quasi-cohérent sur $\mathcal{O}_X$.
Supposons que $f$ soit un morphisme fini.
Alors $\mathcal{F}$ est un module sur $\mathcal{O}_X$ de type fini
si et seulement si $f_*\mathcal{F}$ est un module sur $\mathcal{O}_Y$ de
fini.
\end{lemma}

\begin{proof}
Comme $f$ est fini, il est représentable. Choisissons un schéma $V$ et un
morphisme \'etale surjectif $V \to Y$. Alors $U = V \times_Y X$ est un schéma muni d'un
morphisme \'etale surjectif vers $X$ et d'un morphisme fini
$\psi : U \to V$ (le changement de base de $f$). Puisque
$\psi_*(\mathcal{F}|_U) = f_*\mathcal{F}|_V$
le résultat du lemme découle immédiatement de la version pour les schémas, qui
est
Descente, lemme \ref{descent-lemma-finite-over-finite-module}.
\end{proof}

\begin{lemma}
\label{lemma-finite-finitely-presented-module}
Soit $S$ un schéma.
Soit $f : X \to Y$ un morphisme d'espaces algébriques sur $S$.
Soit $\mathcal{F}$ un module quasi-cohérent sur $\mathcal{O}_X$.
Supposons que $f$ soit fini et de présentation finie.
Alors $\mathcal{F}$ est un module sur $\mathcal{O}_X$ de présentation finie
si et seulement si $f_*\mathcal{F}$ est un module sur $\mathcal{O}_Y$ de
présentation.
\end{lemma}

\begin{proof}
Comme $f$ est fini, il est représentable. Choisissons un schéma $V$ et un
morphisme \'etale surjectif $V \to Y$. Alors $U = V \times_Y X$ est un schéma muni d'un
morphisme \'etale surjectif vers $X$ et d'un morphisme fini
$\psi : U \to V$ (le changement de base de $f$). Puisque
$\psi_*(\mathcal{F}|_U) = f_*\mathcal{F}|_V$
le résultat du lemme découle immédiatement de la version pour les schémas, qui
est
Descente, lemme \ref{descent-lemma-finite-finitely-presented-module}.
\end{proof}






\section{Recouvrements fpqc}
\label{section-fpqc}

\noindent
Cette section est l'analogue de
Descente, section \ref{descent-section-fpqc-universal-effective-epimorphisms}.
Pour l'instant, nous ne savons pas si tous les résultats concernant les
recouvrements fpqc de schémas restent valables pour les espaces algébriques.

\begin{lemma}
\label{lemma-open-fpqc-covering}
Soit $S$ un schéma.
Soit $\{f_i : T_i \to T\}_{i \in I}$ un recouvrement fpqc
d'espaces algébriques sur $S$.
Supposons que, pour chaque $i$, nous ayons un sous-espace ouvert $W_i \subset T_i$
tel que, pour tous $i, j \in I$, nous ayons
$\text{pr}_0^{-1}(W_i) = \text{pr}_1^{-1}(W_j)$ comme sous-espaces ouverts
sous-espaces ouverts de $T_i \times_T T_j$. Alors il existe un unique sous-espace ouvert
$W \subset T$ tel que $W_i = f_i^{-1}(W)$ pour tout $i$.
\end{lemma}

\begin{proof}
Par
Topologies sur les espaces, lemme \ref{spaces-topologies-lemma-refine-fpqc-schemes}
nous pouvons supposer que chaque $T_i$ est un schéma.
Choisissons un schéma $U$ et un morphisme \'etale surjectif $U \to T$.
Alors $\{T_i \times_T U \to U\}$ est un recouvrement fpqc de $U$
et $T_i \times_T U$ est un schéma pour chaque $i$. Ainsi
la famille des ouverts $W_i \times_T U$ provient d'un unique
sous-schéma ouvert $W' \subset U$ par
Descente, lemme \ref{descent-lemma-open-fpqc-covering}.
Comme $U \to X$ est ouvert, nous pouvons définir $W \subset X$ comme l'ouvert de Zariski
qui est l'image de $W'$, voir
Propriétés des espaces, section \ref{spaces-properties-section-points}.
Nous omettons la vérification que cela fonctionne, c'est-à-dire que
$W_i$ est l'image inverse de $W$ pour chaque $i$.
\end{proof}

\begin{lemma}
\label{lemma-fpqc-universal-effective-epimorphisms}
Soit $S$ un schéma. Soit $\{T_i \to T\}$ un recouvrement fpqc d'espaces algébriques
sur $S$, voir Topologies sur les espaces, définition
\ref{spaces-topologies-definition-fpqc-covering}.
Alors, pour tout espace algébrique $B$ sur $S$, la suite
$$
\xymatrix{
\Mor_S(T, B) \ar[r] &
\prod\nolimits_i \Mor_S(T_i, B) \ar@<1ex>[r] \ar@<-1ex>[r] &
\prod\nolimits_{i, j} \Mor_S(T_i \times_T T_j, B)
}
$$
est un diagramme égaliseur.
Autrement dit, tout foncteur représentable sur la catégorie des
espaces algébriques sur $S$ satisfait la condition de faisceau pour les
recouvrements fpqc.
\end{lemma}

\begin{proof}
Nous savons que cela est vrai si $\{T_i \to T\}$ est un recouvrement fpqc de
schémas, voir Propriétés des espaces, proposition
\ref{spaces-properties-proposition-sheaf-fpqc}.
C'est le fait essentiel et nous encourageons le lecteur à passer le reste
de la démonstration, qui est formelle. Choisissons un schéma $U$ et un
morphisme \'etale surjectif
$U \to T$. Soit $U_i$ un schéma et soit $U_i \to T_i \times_T U$
soit un morphisme \'etale surjectif. Alors $\{U_i \to U\}$ est un
recouvrement fpqc. Cela découle de
Topologies sur les espaces, lemmes \ref{spaces-topologies-lemma-fpqc} et
\ref{spaces-topologies-lemma-recognize-fpqc-covering}.
D'après ce qui précède, le résultat vaut pour $\{U_i \to U\}$.

\medskip\noindent
Voici ce que cela signifie : supposons que $b_i : T_i \to B$
soit une famille de morphismes telle que
$b_i \circ \text{pr}_0 = b_j \circ \text{pr}_1$ comme morphismes
$T_i \times_T T_j \to B$. Posons alors $a_i : U_i \to B$ égal au
composé de $U_i \to T_i$ et de $b_i$. D'après ce qui précède
nous trouvons un unique morphisme $a : U \to B$ tel que
$a_i$ est le composé de $a$ avec $U_i \to U$.
L'unicité garantit que $a \circ \text{pr}_0 = a \circ \text{pr}_1$
comme morphismes $U \times_T U \to B$. Puisque $T = U/(U \times_T U)$
comme faisceau, nous trouvons que $a$ provient d'un unique morphisme $b : T \to B$.
En poursuivant les diagrammes, nous trouvons que $b$ est le morphisme recherché.
\end{proof}










\section{Descente des propriétés de finitude et de lissité des morphismes}
\label{section-descent-finiteness-morphisms}

\noindent
Le type de lemme suivant est parfois utile.

\begin{lemma}
\label{lemma-curiosity}
Soit $S$ un schéma. Soient $X \to Y \to Z$ des morphismes d'espaces algébriques.
Soit $P$ l'une des propriétés suivantes des morphismes d'espaces algébriques
sur $S$ :
plate, localement de type fini, localement de présentation finie.
Supposons que $X \to Z$ possède $P$ et que
$X \to Y$ est une surjection de faisceaux sur $(\Sch/S)_{fppf}$.
Alors $Y \to Z$ possède $P$.
\end{lemma}

\begin{proof}
Choisissons un schéma $W$ et un morphisme \'etale surjectif $W \to Z$.
Choisissons un schéma $V$ et un morphisme \'etale surjectif $V \to W \times_Z Y$.
Choisissons un schéma $U$ et un morphisme \'etale surjectif $U \to V \times_Y X$.
Par hypothèse, il existe un recouvrement fppf $\{V_i \to V\}$ et des
relèvements $V_i \to X$ du morphisme $V_i \to Y$. Comme $U \to X$ est surjectif
\'etale, nous voyons qu'au-dessus des membres du recouvrement fppf
$\{V_i \times_X U \to V\}$ nous avons des relèvements vers $U$. Ainsi $U \to V$ induit
une surjection de faisceaux sur $(\Sch/S)_{fppf}$.
Par notre définition de ce que signifie avoir la propriété $P$ pour un
morphisme d'espaces algébriques (voir
Morphismes d'espaces,
définition \ref{spaces-morphisms-definition-flat},
définition \ref{spaces-morphisms-definition-locally-finite-type}, et
définition \ref{spaces-morphisms-definition-locally-finite-presentation})
nous voyons que $U \to W$ possède $P$ et qu'il faut montrer que $V \to W$ possède $P$.
Nous nous ramenons ainsi au cas des morphismes de schémas
traité dans
Descente, lemme \ref{descent-lemma-curiosity}.
\end{proof}

\noindent
Un cas plus classique du lemme précédent est le suivant.
(La version « plate » découle de
Morphismes d'espaces, lemme \ref{spaces-morphisms-lemma-flat-permanence}.)

\begin{lemma}
\label{lemma-flat-finitely-presented-permanence}
Soit $S$ un schéma. Soit
$$
\xymatrix{
X \ar[rr]_f \ar[rd]_p & &
Y \ar[dl]^q \\
& B
}
$$
soit un diagramme commutatif de morphismes d'espaces algébriques sur $S$.
Supposons que $f$ soit surjectif, plat et localement de présentation finie
et que $p$ soit localement de présentation finie (resp.\ localement
de type fini). Alors $q$ est localement de présentation finie
(resp.\ localement de type fini).
\end{lemma}

\begin{proof}
Comme $\{X \to Y\}$ est un recouvrement fppf, il induit une surjection de
faisceaux fppf (Topologies sur les espaces, lemme
\ref{spaces-topologies-lemma-fppf-covering-surjective}) et le
le lemme est un cas particulier du lemme \ref{lemma-curiosity}.
D'autre part, un argument plus simple consiste à le déduire de
l'analogue pour les schémas. En effet, le problème est local pour la topologie \'etale
sur $B$ et $Y$ (Morphismes d'espaces, lemmes
\ref{spaces-morphisms-lemma-finite-type-local} et
\ref{spaces-morphisms-lemma-finite-presentation-local}).
Nous pouvons donc supposer que $B$ et $Y$ sont des
schémas affines. Comme $|X| \to |Y|$ est ouvert
(Morphismes d'espaces, lemme \ref{spaces-morphisms-lemma-fppf-open}),
nous pouvons choisir un
schéma affine $U$ et un morphisme \'etale $U \to X$ tel que le
morphisme composé $U \to Y$ soit surjectif. Dans ce cas, le résultat
découle de Descente, lemme
\ref{descent-lemma-flat-finitely-presented-permanence}.
\end{proof}

\begin{lemma}
\label{lemma-syntomic-smooth-etale-permanence}
Soit $S$ un schéma. Soit
$$
\xymatrix{
X \ar[rr]_f \ar[rd]_p & &
Y \ar[dl]^q \\
& B
}
$$
soit un diagramme commutatif de morphismes d'espaces algébriques sur $S$.
Supposons que
\begin{enumerate}
\item $f$ soit surjectif et syntomique (resp.\ lisse, resp.\ \'etale),
\item $p$ soit syntomique (resp.\ lisse, resp.\ \'etale).
\end{enumerate}
Alors $q$ est syntomique (resp.\ lisse, resp.\ \'etale).
\end{lemma}

\begin{proof}
Nous le déduisons de l'analogue pour les schémas.
En effet, le problème est local pour la topologie \'etale sur $B$ et $Y$
(Morphismes d'espaces, lemmes
\ref{spaces-morphisms-lemma-syntomic-local},
\ref{spaces-morphisms-lemma-smooth-local}, et
\ref{spaces-morphisms-lemma-etale-local}).
Nous pouvons donc supposer que $B$ et $Y$ sont des
schémas affines. Comme $|X| \to |Y|$ est ouvert
(Morphismes d'espaces, lemme \ref{spaces-morphisms-lemma-fppf-open}),
nous pouvons choisir un
schéma affine $U$ et un morphisme \'etale $U \to X$ tel que le
morphisme composé $U \to Y$ soit surjectif. Dans ce cas, le résultat
découle de Descente, lemme
\ref{descent-lemma-syntomic-smooth-etale-permanence}.
\end{proof}

\noindent
En fait, nous pouvons renforcer ce résultat comme suit.

\begin{lemma}
\label{lemma-smooth-permanence}
Soit $S$ un schéma. Soit
$$
\xymatrix{
X \ar[rr]_f \ar[rd]_p & &
Y \ar[dl]^q \\
& B
}
$$
soit un diagramme commutatif de morphismes d'espaces algébriques sur $S$. Supposons que
\begin{enumerate}
\item $f$ soit surjectif, plat et localement de présentation finie,
\item $p$ soit lisse (resp.\ \'etale).
\end{enumerate}
Alors $q$ est lisse (resp.\ \'etale).
\end{lemma}

\begin{proof}
Nous le déduisons de l'analogue pour les schémas.
En effet, le problème est local pour la topologie \'etale sur $B$ et $Y$
(Morphismes d'espaces, lemmes
\ref{spaces-morphisms-lemma-smooth-local} et
\ref{spaces-morphisms-lemma-etale-local}).
Nous pouvons donc supposer que $B$ et $Y$ sont des
schémas affines. Comme $|X| \to |Y|$ est ouvert
(Morphismes d'espaces, lemme \ref{spaces-morphisms-lemma-fppf-open}),
nous pouvons choisir un
schéma affine $U$ et un morphisme \'etale $U \to X$ tel que le
morphisme composé $U \to Y$ soit surjectif. Dans ce cas, le résultat
découle de Descente, lemme
\ref{descent-lemma-smooth-permanence}.
\end{proof}

\begin{lemma}
\label{lemma-syntomic-permanence}
Soit $S$ un schéma. Soit
$$
\xymatrix{
X \ar[rr]_f \ar[rd]_p & &
Y \ar[dl]^q \\
& B
}
$$
soit un diagramme commutatif de morphismes d'espaces algébriques sur $S$. Supposons que
\begin{enumerate}
\item $f$ soit surjectif, plat et localement de présentation finie,
\item $p$ soit syntomique.
\end{enumerate}
Alors $q$ et $f$ sont tous deux syntomiques.
\end{lemma}

\begin{proof}
Nous le déduisons de l'analogue pour les schémas.
En effet, le problème est local pour la topologie \'etale sur $B$ et $Y$
(Morphismes d'espaces, lemme
\ref{spaces-morphisms-lemma-syntomic-local}).
Nous pouvons donc supposer que $B$ et $Y$ sont des
schémas affines. Comme $|X| \to |Y|$ est ouvert
(Morphismes d'espaces, lemme \ref{spaces-morphisms-lemma-fppf-open}),
nous pouvons choisir un
schéma affine $U$ et un morphisme \'etale $U \to X$ tel que le
morphisme composé $U \to Y$ soit surjectif. Dans ce cas, le résultat
découle de Descente, lemme
\ref{descent-lemma-syntomic-permanence}.
\end{proof}










\section{Descente de propriétés des espaces}
\label{section-descending-properties-spaces}

\noindent
Dans cette section, nous rassemblons quelques résultats du type suivant.

\begin{lemma}
\label{lemma-descend-unibranch}
Soit $S$ un schéma.
Soit $f : X \to Y$ un morphisme d'espaces algébriques sur $S$.
Soit $x \in |X|$.
Si $f$ est plat en $x$ et si $X$ est géométriquement unibranche en $x$, alors $Y$ est
géométriquement unibranche en $f(x)$.
\end{lemma}

\begin{proof}
Considérons le morphisme d'anneaux locaux \'etale
$\mathcal{O}_{Y, f(\overline{x})} \to \mathcal{O}_{X, \overline{x}}$.
Par
Morphismes d'espaces, lemme
\ref{spaces-morphisms-lemma-flat-at-point-etale-local-rings}
celui-ci est plat. Ainsi, si $\mathcal{O}_{X, \overline{x}}$ possède un unique idéal premier minimal
il en va de même de $\mathcal{O}_{Y, f(\overline{x})}$ (par descente, voir
Algèbre, lemme \ref{algebra-lemma-flat-going-down}).
\end{proof}

\begin{lemma}
\label{lemma-descend-reduced}
\begin{slogan}
Un morphisme plat et surjectif d'espaces algébriques dont la source est réduite
a un but réduit.
\end{slogan}
Soit $S$ un schéma.
Soit $f : X \to Y$ un morphisme d'espaces algébriques sur $S$.
Si $f$ est plat et surjectif et si $X$ est réduit, alors $Y$ est réduit.
\end{lemma}

\begin{proof}
Choisissons un schéma $V$ et un morphisme \'etale surjectif $V \to Y$.
Choisissons un schéma $U$ et un morphisme \'etale surjectif
$U \to X \times_Y V$. Comme $f$ est surjectif et plat, le morphisme de
schémas $U \to V$ est surjectif et plat. Nous nous ramenons ainsi
au cas des schémas (la réduction de $X$ et $Y$ étant définie
en termes de réduction de $U$ et $V$, voir
Propriétés des espaces,
section \ref{spaces-properties-section-types-properties}).
Le cas des schémas est traité dans
Descente, lemme \ref{descent-lemma-descend-reduced}.
\end{proof}

\begin{lemma}
\label{lemma-descend-locally-Noetherian}
Soit $f : X \to Y$ un morphisme d'espaces algébriques.
Si $f$ est localement de présentation finie, plat et surjectif et si
$X$ est localement noethérien, alors $Y$ est localement noethérien.
\end{lemma}

\begin{proof}
Choisissons un schéma $V$ et un morphisme \'etale surjectif $V \to Y$.
Choisissons un schéma $U$ et un morphisme \'etale surjectif
$U \to X \times_Y V$. Comme $f$ est surjectif, plat et localement de
présentation finie, le morphisme de schémas $U \to V$ est surjectif, plat et
localement de présentation finie. Nous nous ramenons ainsi
au cas des schémas (la propriété d'être localement noethérien pour $X$ et $Y$
étant définie en termes de la propriété d'être localement noethérien de $U$ et $V$, voir
Propriétés des espaces,
section \ref{spaces-properties-section-types-properties}).
Dans le cas des schémas, le résultat découle de
Descente, lemme \ref{descent-lemma-Noetherian-local-fppf}.
\end{proof}

\begin{lemma}
\label{lemma-descend-regular}
Soit $f : X \to Y$ un morphisme d'espaces algébriques.
Si $f$ est localement de présentation finie, plat et surjectif et si
$X$ est régulier, alors $Y$ est régulier.
\end{lemma}

\begin{proof}
Par
Lemme \ref{lemma-descend-locally-Noetherian}
nous savons que $Y$ est localement noethérien.
Choisissons un schéma $V$ et un morphisme \'etale surjectif $V \to Y$.
Il suffit de montrer que les anneaux locaux de $V$ sont tous des anneaux locaux réguliers
, voir
Propriétés, lemme \ref{properties-lemma-characterize-regular}.
Choisissons un schéma $U$ et un morphisme \'etale surjectif
$U \to X \times_Y V$. Comme $f$ est surjectif et plat, le morphisme de schémas
$U \to V$ est surjectif et plat. Par hypothèse, $U$ est un schéma régulier
en particulier tous ses anneaux locaux sont réguliers (par le lemme ci-dessus).
Le lemme découle donc de
Algèbre, lemme \ref{algebra-lemma-flat-under-regular}.
\end{proof}

\begin{lemma}
\label{lemma-reduced-local-smooth}
Soit $f : X \to Y$ un morphisme lisse d'espaces algébriques.
Si $Y$ est réduit, alors $X$ est réduit. Si $f$ est surjectif
et si $X$ est réduit, alors $Y$ est réduit.
\end{lemma}

\begin{proof}
Choisissons un diagramme commutatif
$$
\xymatrix{
U \ar[d] \ar[r] & V \ar[d] \\
X \ar[r] & Y
}
$$
où $U$ et $V$ sont des schémas, les flèches verticales sont surjectives et
\'etales, et $U \to X \times_Y V$ est \'etale surjectif. Observons que $X$ est
un espace algébrique réduit si et seulement si $U$ est un schéma réduit
par notre définition des espaces algébriques réduits dans
Propriétés des espaces, section \ref{spaces-properties-section-types-properties}.
De même pour $Y$ et $V$. 
Le morphisme $U \to V$ est un morphisme lisse de schémas, voir
Morphismes d'espaces, lemme \ref{spaces-morphisms-lemma-smooth-local}.
Puisque le fait d'être réduit est local pour la topologie lisse pour les
schémas (Descente, lemme \ref{descent-lemma-reduced-local-smooth})
nous voyons que $U$ est réduit si $V$ est réduit.
D'autre part, si $X \to Y$ est surjectif, alors $U \to V$ est
surjectif et, dans ce cas, si $U$ est réduit, alors $V$ est réduit.
\end{proof}







\section{Descente de propriétés des morphismes}
\label{section-descending-properties-morphisms}

\noindent
Dans cette section, nous introduisons la notion selon laquelle une propriété des morphismes d'
espaces algébriques est locale sur le but pour une topologie. Comparer avec
Descente, section \ref{descent-section-descending-properties-morphisms}.

\begin{definition}
\label{definition-property-morphisms-local}
Soit $S$ un schéma.
Soit $\mathcal{P}$ une propriété des morphismes d'espaces algébriques sur $S$.
Soit $\tau \in \{fpqc, fppf, syntomic, smooth, \etale\}$.
On dit que $\mathcal{P}$ est {\it locale pour $\tau$ sur la base}, ou
{\it locale pour $\tau$ sur le but}, ou
{\it locale sur la base pour la topologie-$\tau$} si, pour tout
recouvrement-$\tau$ $\{Y_i \to Y\}_{i \in I}$ d'espaces algébriques
et tout morphisme d'espaces algébriques $f : X \to Y$, nous
have
$$
f \text{ has }\mathcal{P}
\Leftrightarrow
\text{each }Y_i \times_Y X \to Y_i\text{ has }\mathcal{P}.
$$
\end{definition}

\noindent
To be sure, since isomorphisms are always coverings
we see (or require) that property $\mathcal{P}$ holds for $X \to Y$
if and only if it holds for any arrow $X' \to Y'$ isomorphic to $X \to Y$.
If a property is $\tau$-local on the target then it is preserved
by base changes by morphisms which occur in $\tau$-coverings. Here
is a formal statement.

\begin{lemma}
\label{lemma-pullback-property-local-target}
Let $S$ be a scheme.
Let $\tau \in \{fpqc, fppf, syntomic, smooth, \etale\}$.
Let $\mathcal{P}$ be a property of morphisms of algebraic spaces over $S$
which is $\tau$ local on the target. Let $f : X \to Y$ have property
$\mathcal{P}$. For any morphism $Y' \to Y$ which is
flat, resp.\ flat and locally of finite presentation, resp.\ syntomic,
resp.\ \'etale, the base change
$f' : Y' \times_Y X \to Y'$ of $f$ has property $\mathcal{P}$.
\end{lemma}

\begin{proof}
This is true because we can fit $Y' \to Y$ into a family of
morphisms which forms a $\tau$-covering.
\end{proof}

\noindent
A simple often used consequence of the above is that if
$f : X \to Y$ has property $\mathcal{P}$ which is $\tau$-local
on the target and $f(X) \subset V$
for some open subspace $V \subset Y$, then also the induced
morphism $X \to V$ has $\mathcal{P}$. Proof: The base change
$f$ by $V \to Y$ gives $X \to V$.

\begin{lemma}
\label{lemma-largest-open-of-the-base}
Let $S$ be a scheme.
Let $\tau \in \{fppf, syntomic, smooth, \etale\}$.
Let $\mathcal{P}$ be a property of morphisms of algebraic spaces over $S$
which is $\tau$ local on the target. For any morphism of algebraic spaces
$f : X \to Y$ over $S$ there exists a largest open subspace
$W(f) \subset Y$ such that the restriction $X_{W(f)} \to W(f)$ has
$\mathcal{P}$. Moreover,
\begin{enumerate}
\item if $g : Y' \to Y$ is a morphism of algebraic spaces which is
flat and locally of finite presentation, syntomic, smooth, or \'etale
and the base change $f' : X_{Y'} \to Y'$ has $\mathcal{P}$, then
$g$ factors through $W(f)$,
\item if $g : Y' \to Y$ is flat and locally of finite presentation,
syntomic, smooth, or \'etale, then $W(f') = g^{-1}(W(f))$, and
\item if $\{g_i : Y_i \to Y\}$ is a $\tau$-covering, then
$g_i^{-1}(W(f)) = W(f_i)$, where $f_i$ is the base change of $f$
by $Y_i \to Y$.
\end{enumerate}
\end{lemma}

\begin{proof}
Consider the union $W_{set} \subset |Y|$ of the images
$g(|Y'|) \subset |Y|$ of morphisms $g : Y' \to Y$ with the properties:
\begin{enumerate}
\item $g$ is flat and locally of finite presentation, syntomic,
smooth, or \'etale, and
\item the base change $Y' \times_{g, Y} X \to Y'$ has property
$\mathcal{P}$.
\end{enumerate}
Since such a morphism $g$ is open (see
Morphisms of Spaces, Lemma \ref{spaces-morphisms-lemma-fppf-open})
we see that $W_{set}$ is an open subset of $|Y|$. Denote $W \subset Y$
the open subspace whose underlying set of points is $W_{set}$, see
Properties of Spaces, Lemma \ref{spaces-properties-lemma-open-subspaces}.
Since $\mathcal{P}$ is local in the $\tau$ topology the restriction
$X_W \to W$ has property $\mathcal{P}$ because we are given a covering
$\{Y' \to W\}$ of $W$ such that the pullbacks have $\mathcal{P}$.
This proves the existence and proves that $W(f)$ has property (1).
To see property (2) note that $W(f') \supset g^{-1}(W(f))$ because
$\mathcal{P}$ is stable under base change by flat and locally of finite
presentation, syntomic, smooth, or \'etale morphisms, see
Lemma \ref{lemma-pullback-property-local-target}.
On the other hand, if $Y'' \subset Y'$ is an open such that
$X_{Y''} \to Y''$ has property $\mathcal{P}$, then $Y'' \to Y$ factors
through $W$ by construction, i.e., $Y'' \subset g^{-1}(W(f))$. This
proves (2). Assertion (3) follows from (2) because each morphism
$Y_i \to Y$ is flat and locally of finite presentation, syntomic,
smooth, or \'etale by our definition of a $\tau$-covering.
\end{proof}

\begin{lemma}
\label{lemma-descending-properties-morphisms}
Let $S$ be a scheme. Let $\mathcal{P}$ be a property of morphisms of
algebraic spaces over $S$. Assume
\begin{enumerate}
\item if $X_i \to Y_i$, $i = 1, 2$ have property $\mathcal{P}$ so
does $X_1 \amalg X_2 \to Y_1 \amalg Y_2$,
\item a morphism of algebraic spaces $f : X \to Y$ has property
$\mathcal{P}$ if and only if for every affine scheme $Z$ and
morphism $Z \to Y$ the base change $Z \times_Y X \to Z$ of $f$
has property $\mathcal{P}$, and
\item for any surjective flat morphism of affine schemes
$Z' \to Z$ over $S$ and a morphism $f : X \to Z$ from an algebraic space
to $Z$ we have
$$
f' : Z' \times_Z X \to Z'\text{ has }\mathcal{P}
\Rightarrow
f\text{ has }\mathcal{P}.
$$
\end{enumerate}
Then $\mathcal{P}$ is fpqc local on the base.
\end{lemma}

\begin{proof}
If $\mathcal{P}$ has property (2), then it is automatically
stable under any base change. Hence the direct implication in
Definition \ref{definition-property-morphisms-local}.

\medskip\noindent
Let $\{Y_i \to Y\}_{i \in I}$ be an fpqc covering of algebraic spaces over $S$.
Let $f : X \to Y$ be a morphism of algebraic spaces over $S$.
Assume each base change $f_i : Y_i \times_Y X \to Y_i$ has property
$\mathcal{P}$. Our goal is to show that $f$ has $\mathcal{P}$.
Let $Z$ be an affine scheme, and let $Z \to Y$ be a morphism.
By (2) it suffices to show that the morphism of algebraic spaces
$Z \times_Y X \to Z$ has $\mathcal{P}$.
Since $\{Y_i \to Y\}_{i \in I}$ is an fpqc covering we know there
exists a standard fpqc covering $\{Z_j \to Z\}_{j = 1, \ldots , n}$
and morphisms $Z_j \to Y_{i_j}$ over $Y$ for suitable indices $i_j \in I$.
Since $f_{i_j}$ has $\mathcal{P}$ we see that
$$
Z_j \times_Y X
=
Z_j \times_{Y_{i_j}} (Y_{i_j} \times_Y X)
\longrightarrow
Z_j
$$
has $\mathcal{P}$ as a base change of $f_{i_j}$ (see first remark of the
proof). Set $Z' = \coprod_{j = 1, \ldots, n} Z_j$, so that $Z' \to Z$ is
a flat and surjective morphism of affine schemes over $S$. By (1)
we conclude that $Z' \times_Y X \to Z'$ has property $\mathcal{P}$.
Since this is the base change of the morphism $Z \times_Y X \to Z$
by the morphism $Z' \to Z$ we conclude that
$Z \times_Y X \to Z$ has property $\mathcal{P}$ as desired.
\end{proof}


\section{Descending properties of morphisms in the fpqc topology}
\label{section-descending-properties-morphisms-fpqc}


\noindent
In this section we find a large number of properties
of morphisms of algebraic spaces which are local on the base
in the fpqc topology. Please compare with
Descent, Section \ref{descent-section-descending-properties-morphisms-fpqc}
for the case of morphisms of schemes.

\begin{lemma}
\label{lemma-descending-property-quasi-compact}
Let $S$ be a scheme.
The property $\mathcal{P}(f) =$``$f$ is quasi-compact''
is fpqc local on the base on algebraic spaces over $S$.
\end{lemma}

\begin{proof}
We will use
Lemma \ref{lemma-descending-properties-morphisms}
to prove this. Assumptions (1) and (2) of that lemma follow from
Morphisms of Spaces,
Lemma \ref{spaces-morphisms-lemma-quasi-compact-local}.
Let $Z' \to Z$ be a surjective flat morphism of affine schemes over $S$.
Let $f : X \to Z$ be a morphism of algebraic spaces, and assume
that the base change $f' : Z' \times_Z X \to Z'$ is quasi-compact. We have
to show that $f$ is quasi-compact. To see this, using
Morphisms of Spaces,
Lemma \ref{spaces-morphisms-lemma-quasi-compact-local}
again, it is enough to show that for every affine scheme $Y$ and
morphism $Y \to Z$ the fibre product $Y \times_Z X$ is quasi-compact.
Here is a picture:
\begin{equation}
\label{equation-cube}
\vcenter{
\xymatrix{
Y \times_Z Z' \times_Z X \ar[dd] \ar[rr] \ar[rd] & &
Z' \times_Z X \ar'[d][dd]^{f'} \ar[rd] \\
& Y \times_Z X \ar[dd] \ar[rr] & & X \ar[dd]^f \\
Y \times_Z Z' \ar'[r][rr] \ar[rd] & & Z' \ar[rd] \\
& Y \ar[rr] & & Z
}
}
\end{equation}
Note that all squares are cartesian and the bottom square consists
of affine schemes. The assumption that $f'$ is quasi-compact combined with
the fact that $Y \times_Z Z'$ is affine implies that
$Y \times_Z Z' \times_Z X$ is quasi-compact. Since
$$
Y \times_Z Z' \times_Z X \longrightarrow Y \times_Z X
$$
is surjective as a base change of $Z' \to Z$
we conclude that $Y \times_Z X$ is quasi-compact, see
Morphisms of Spaces,
Lemma \ref{spaces-morphisms-lemma-surjection-from-quasi-compact}.
This finishes the proof.
\end{proof}

\begin{lemma}
\label{lemma-descending-property-quasi-separated}
Let $S$ be a scheme.
The property $\mathcal{P}(f) =$``$f$ is quasi-separated''
is fpqc local on the base on algebraic spaces over $S$.
\end{lemma}

\begin{proof}
A base change of a quasi-separated morphism is quasi-separated, see
Morphisms of Spaces,
Lemma \ref{spaces-morphisms-lemma-base-change-separated}.
Hence the direct implication in
Definition \ref{definition-property-morphisms-local}.

\medskip\noindent
Let $\{Y_i \to Y\}_{i \in I}$ be an fpqc covering of algebraic spaces over $S$.
Let $f : X \to Y$ be a morphism of algebraic spaces over $S$.
Assume each base change $X_i := Y_i \times_Y X \to Y_i$ is quasi-separated.
This means that each of the morphisms
$$
\Delta_i :
X_i
\longrightarrow
X_i \times_{Y_i} X_i = Y_i \times_Y (X \times_Y X)
$$
is quasi-compact. The base change of a fpqc covering is an fpqc covering, see
Topologies on Spaces, Lemma \ref{spaces-topologies-lemma-fpqc}
hence $\{Y_i \times_Y (X \times_Y X) \to X \times_Y X\}$
is an fpqc covering of algebraic spaces. Moreover, each
$\Delta_i$ is the base change of the morphism
$\Delta : X \to X \times_Y X$. Hence it follows from
Lemma \ref{lemma-descending-property-quasi-compact}
that $\Delta$ is quasi-compact, i.e., $f$ is quasi-separated.
\end{proof}

\begin{lemma}
\label{lemma-descending-property-universally-closed}
Let $S$ be a scheme.
The property $\mathcal{P}(f) =$``$f$ is universally closed''
is fpqc local on the base on algebraic spaces over $S$.
\end{lemma}

\begin{proof}
We will use
Lemma \ref{lemma-descending-properties-morphisms}
to prove this. Assumptions (1) and (2) of that lemma follow from
Morphisms of Spaces,
Lemma \ref{spaces-morphisms-lemma-universally-closed-local}.
Let $Z' \to Z$ be a surjective flat morphism of affine schemes over $S$.
Let $f : X \to Z$ be a morphism of algebraic spaces, and assume
that the base change $f' : Z' \times_Z X \to Z'$ is universally closed.
We have to show that $f$ is universally closed. To see this, using
Morphisms of Spaces,
Lemma \ref{spaces-morphisms-lemma-universally-closed-local}
again, it is enough to show that for every affine scheme $Y$ and
morphism $Y \to Z$ the map $|Y \times_Z X| \to |Y|$ is closed.
Consider the cube (\ref{equation-cube}).
The assumption that $f'$ is universally closed implies that
$|Y \times_Z Z' \times_Z X| \to |Y \times_Z Z'|$ is closed.
As $Y \times_Z Z' \to Y$ is quasi-compact, surjective, and flat
as a base change of $Z' \to Z$
we see the map $|Y \times_Z Z'| \to |Y|$ is submersive, see
Morphisms, Lemma \ref{morphisms-lemma-fpqc-quotient-topology}.
Moreover the map
$$
|Y \times_Z Z' \times_Z X|
\longrightarrow
|Y \times_Z Z'| \times_{|Y|} |Y \times_Z X|
$$
is surjective, see
Properties of Spaces, Lemma \ref{spaces-properties-lemma-points-cartesian}.
It follows by elementary topology that $|Y \times_Z X| \to |Y|$ is closed.
\end{proof}

\begin{lemma}
\label{lemma-descending-property-universally-open}
Let $S$ be a scheme.
The property $\mathcal{P}(f) =$``$f$ is universally open''
is fpqc local on the base on algebraic spaces over $S$.
\end{lemma}

\begin{proof}
The proof is the same as the proof of
Lemma \ref{lemma-descending-property-universally-closed}.
\end{proof}

\begin{lemma}
\label{lemma-descending-property-universally-submersive}
The property $\mathcal{P}(f) =$``$f$ is universally submersive''
is fpqc local on the base.
\end{lemma}

\begin{proof}
The proof is the same as the proof of
Lemma \ref{lemma-descending-property-universally-closed}.
\end{proof}

\begin{lemma}
\label{lemma-descending-property-surjective}
The property $\mathcal{P}(f) =$``$f$ is surjective''
is fpqc local on the base.
\end{lemma}

\begin{proof}
Omitted. (Hint: Use
Properties of Spaces, Lemma \ref{spaces-properties-lemma-points-cartesian}.)
\end{proof}

\begin{lemma}
\label{lemma-descending-property-universally-injective}
The property $\mathcal{P}(f) =$``$f$ is universally injective''
is fpqc local on the base.
\end{lemma}

\begin{proof}
We will use
Lemma \ref{lemma-descending-properties-morphisms}
to prove this. Assumptions (1) and (2) of that lemma follow from
Morphisms of Spaces,
Lemma \ref{spaces-morphisms-lemma-universally-closed-local}.
Let $Z' \to Z$ be a flat surjective morphism of affine schemes
over $S$ and let $f : X \to Z$ be a morphism from an algebraic space to $Z$.
Assume that the base change $f' : X' \to Z'$ is universally injective.
Let $K$ be a field, and let $a, b : \Spec(K) \to X$
be two morphisms such that $f \circ a = f \circ b$.
As $Z' \to Z$ is surjective there exists a field
extension $K'/K$ and a morphism
$\Spec(K') \to Z'$
such that the following solid diagram commutes
$$
\xymatrix{
\Spec(K') \ar[rrd] \ar@{-->}[rd]_{a', b'} \ar[dd] \\
 &
X' \ar[r] \ar[d] &
Z' \ar[d] \\
\Spec(K) \ar[r]^{a, b} &
X \ar[r] &
Z
}
$$
As the square is cartesian we get the two dotted arrows $a'$, $b'$ making the
diagram commute. Since $X' \to Z'$ is universally injective we get $a' = b'$.
This forces $a = b$ as $\{\Spec(K') \to \Spec(K)\}$
is an fpqc covering, see
Properties of Spaces, Proposition
\ref{spaces-properties-proposition-sheaf-fpqc}.
Hence $f$ is universally injective as desired.
\end{proof}

\begin{lemma}
\label{lemma-descending-property-universal-homeomorphism}
The property $\mathcal{P}(f) =$``$f$ is a universal homeomorphism''
is fpqc local on the base.
\end{lemma}

\begin{proof}
This can be proved in exactly the same manner as
Lemma \ref{lemma-descending-property-universally-closed}.
Alternatively, one can use that a map of topological spaces is a
homeomorphism if and only if it is injective, surjective, and open.
Thus a universal homeomorphism is the same thing as a
surjective, universally injective, and universally open morphism.
See Morphisms of Spaces, Lemma
\ref{spaces-morphisms-lemma-base-change-surjective} and
Morphisms of Spaces, Definitions
\ref{spaces-morphisms-definition-universally-injective},
\ref{spaces-morphisms-definition-surjective},
\ref{spaces-morphisms-definition-open},
\ref{spaces-morphisms-definition-universal-homeomorphism}.
Thus the lemma follows from
Lemmas \ref{lemma-descending-property-surjective},
\ref{lemma-descending-property-universally-injective}, and
\ref{lemma-descending-property-universally-open}.
\end{proof}

\begin{lemma}
\label{lemma-descending-property-locally-finite-type}
The property $\mathcal{P}(f) =$``$f$ is locally of finite type''
is fpqc local on the base.
\end{lemma}

\begin{proof}
We will use
Lemma \ref{lemma-descending-properties-morphisms}
to prove this. Assumptions (1) and (2) of that lemma follow from
Morphisms of Spaces,
Lemma \ref{spaces-morphisms-lemma-finite-type-local}.
Let $Z' \to Z$ be a surjective flat morphism of affine schemes over $S$.
Let $f : X \to Z$ be a morphism of algebraic spaces, and assume
that the base change $f' : Z' \times_Z X \to Z'$ is locally of finite type.
We have to show that $f$ is locally of finite type. Let $U$ be a scheme
and let $U \to X$ be surjective and \'etale. By
Morphisms of Spaces,
Lemma \ref{spaces-morphisms-lemma-finite-type-local}
again, it is enough to show that $U \to Z$ is locally of finite type.
Since $f'$ is locally of finite type, and since $Z' \times_Z U$ is a
scheme \'etale over $Z' \times_Z X$ we conclude (by the same lemma again) that
$Z' \times_Z U \to Z'$ is locally of finite type.
As $\{Z' \to Z\}$ is an fpqc covering we conclude that
$U \to Z$ is locally of finite type by
Descent, Lemma \ref{descent-lemma-descending-property-locally-finite-type}
as desired.
\end{proof}

\begin{lemma}
\label{lemma-descending-property-locally-finite-presentation}
The property $\mathcal{P}(f) =$``$f$ is locally of finite presentation''
is fpqc local on the base.
\end{lemma}

\begin{proof}
We will use
Lemma \ref{lemma-descending-properties-morphisms}
to prove this. Assumptions (1) and (2) of that lemma follow from
Morphisms of Spaces,
Lemma \ref{spaces-morphisms-lemma-finite-presentation-local}.
Let $Z' \to Z$ be a surjective flat morphism of affine schemes over $S$.
Let $f : X \to Z$ be a morphism of algebraic spaces, and assume
that the base change $f' : Z' \times_Z X \to Z'$ is locally of
finite presentation.
We have to show that $f$ is locally of finite presentation. Let $U$ be a scheme
and let $U \to X$ be surjective and \'etale. By
Morphisms of Spaces,
Lemma \ref{spaces-morphisms-lemma-finite-presentation-local}
again, it is enough to show that $U \to Z$ is locally of finite presentation.
Since $f'$ is locally of finite presentation, and since $Z' \times_Z U$ is a
scheme \'etale over $Z' \times_Z X$ we conclude (by the same lemma again) that
$Z' \times_Z U \to Z'$ is locally of finite presentation.
As $\{Z' \to Z\}$ is an fpqc covering we conclude that
$U \to Z$ is locally of finite presentation by
Descent,
Lemma \ref{descent-lemma-descending-property-locally-finite-presentation}
as desired.
\end{proof}

\begin{lemma}
\label{lemma-descending-property-finite-type}
The property $\mathcal{P}(f) =$``$f$ is of finite type''
is fpqc local on the base.
\end{lemma}

\begin{proof}
Combine Lemmas \ref{lemma-descending-property-quasi-compact}
and \ref{lemma-descending-property-locally-finite-type}.
\end{proof}

\begin{lemma}
\label{lemma-descending-property-finite-presentation}
The property $\mathcal{P}(f) =$``$f$ is of finite presentation''
is fpqc local on the base.
\end{lemma}

\begin{proof}
Combine Lemmas \ref{lemma-descending-property-quasi-compact},
\ref{lemma-descending-property-quasi-separated} and
\ref{lemma-descending-property-locally-finite-presentation}.
\end{proof}

\begin{lemma}
\label{lemma-descending-property-flat}
The property $\mathcal{P}(f) =$``$f$ is flat''
is fpqc local on the base.
\end{lemma}

\begin{proof}
We will use
Lemma \ref{lemma-descending-properties-morphisms}
to prove this. Assumptions (1) and (2) of that lemma follow from
Morphisms of Spaces,
Lemma \ref{spaces-morphisms-lemma-flat-local}.
Let $Z' \to Z$ be a surjective flat morphism of affine schemes over $S$.
Let $f : X \to Z$ be a morphism of algebraic spaces, and assume
that the base change $f' : Z' \times_Z X \to Z'$ is flat.
We have to show that $f$ is flat. Let $U$ be a scheme
and let $U \to X$ be surjective and \'etale. By
Morphisms of Spaces,
Lemma \ref{spaces-morphisms-lemma-flat-local}
again, it is enough to show that $U \to Z$ is flat.
Since $f'$ is flat, and since $Z' \times_Z U$ is a
scheme \'etale over $Z' \times_Z X$ we conclude (by the same lemma again) that
$Z' \times_Z U \to Z'$ is flat.
As $\{Z' \to Z\}$ is an fpqc covering we conclude that
$U \to Z$ is flat by
Descent, Lemma \ref{descent-lemma-descending-property-flat}
as desired.
\end{proof}

\begin{lemma}
\label{lemma-descending-property-open-immersion}
The property $\mathcal{P}(f) =$``$f$ is an open immersion''
is fpqc local on the base.
\end{lemma}

\begin{proof}
We will use
Lemma \ref{lemma-descending-properties-morphisms}
to prove this. Assumptions (1) and (2) of that lemma follow from
Morphisms of Spaces,
Lemma \ref{spaces-morphisms-lemma-closed-immersion-local}.
Consider a cartesian diagram
$$
\xymatrix{
X' \ar[r] \ar[d] & X \ar[d] \\
Z' \ar[r] & Z
}
$$
of algebraic spaces over $S$
where $Z' \to Z$ is a surjective flat morphism of affine schemes,
and $X' \to Z'$ is an open immersion. We have to show that $X \to Z$
is an open immersion. Note that $|X'| \subset |Z'|$ corresponds to an
open subscheme $U' \subset Z'$ (isomorphic to $X'$)
with the property that $\text{pr}_0^{-1}(U') = \text{pr}_1^{-1}(U')$
as open subschemes of $Z' \times_Z Z'$. Hence there exists an open
subscheme $U \subset Z$ such that $X' = (Z' \to Z)^{-1}(U)$, see
Descent, Lemma \ref{descent-lemma-open-fpqc-covering}.
By Properties of Spaces,
Proposition \ref{spaces-properties-proposition-sheaf-fpqc}
we see that $X$ satisfies the sheaf condition for the fpqc topology.
Now we have the fpqc covering $\mathcal{U} = \{U' \to U\}$
and the element $U' \to X' \to X \in \check{H}^0(\mathcal{U}, X)$.
By the sheaf condition we obtain a morphism $U \to X$ such that
$$
\xymatrix{
U' \ar[r] \ar[d]^{\cong} \ar@/_3ex/[dd] & U \ar[d] \ar@/^3ex/[dd] \\
X' \ar[r] \ar[d] & X \ar[d] \\
Z' \ar[r] & Z
}
$$
is commutative. On the other hand, we know that for any scheme $T$ over $S$
and $T$-valued point $T \to X$ the composition $T \to X \to Z$ is a
morphism such that $Z' \times_Z T \to Z'$ factors through $U'$. Clearly
this means that $T \to Z$ factors through $U$. In other words the map
of sheaves $U \to X$ is bijective and we win.
\end{proof}

\begin{lemma}
\label{lemma-descending-property-isomorphism}
The property $\mathcal{P}(f) =$``$f$ is an isomorphism''
is fpqc local on the base.
\end{lemma}

\begin{proof}
Combine Lemmas \ref{lemma-descending-property-surjective}
and \ref{lemma-descending-property-open-immersion}.
\end{proof}

\begin{lemma}
\label{lemma-descending-property-affine}
The property $\mathcal{P}(f) =$``$f$ is affine''
is fpqc local on the base.
\end{lemma}

\begin{proof}
We will use
Lemma \ref{lemma-descending-properties-morphisms}
to prove this. Assumptions (1) and (2) of that lemma follow from
Morphisms of Spaces,
Lemma \ref{spaces-morphisms-lemma-affine-local}.
Let $Z' \to Z$ be a surjective flat morphism of affine schemes over $S$.
Let $f : X \to Z$ be a morphism of algebraic spaces, and assume
that the base change $f' : Z' \times_Z X \to Z'$ is affine.
Let $X'$ be a scheme representing $Z' \times_Z X$.
We obtain a canonical isomorphism
$$
\varphi : X' \times_Z Z' \longrightarrow Z' \times_Z X'
$$
since both schemes represent the algebraic space $Z' \times_Z Z' \times_Z X$.
This is a descent datum for $X'/Z'/Z$, see
Descent, Definition \ref{descent-definition-descent-datum}
(verification omitted, compare with
Descent, Lemma \ref{descent-lemma-descent-data-sheaves}).
Since $X' \to Z'$ is affine this descent datum is effective, see
Descent, Lemma \ref{descent-lemma-affine}.
Thus there exists a scheme $Y \to Z$ over $Z$ and an
isomorphism $\psi : Z' \times_Z Y \to X'$ compatible with descent data.
Of course $Y \to Z$ is affine (by construction or by
Descent, Lemma \ref{descent-lemma-descending-property-affine}).
Note that $\mathcal{Y} = \{Z' \times_Z Y \to Y\}$ is a
fpqc covering, and interpreting $\psi$ as an element of
$X(Z' \times_Z Y)$ we see that $\psi \in \check{H}^0(\mathcal{Y}, X)$.
By the sheaf condition for $X$ with respect to this covering (see
Properties of Spaces, Proposition
\ref{spaces-properties-proposition-sheaf-fpqc})
we obtain a morphism $Y \to X$.
By construction the base change of this to $Z'$ is an isomorphism, hence
an isomorphism by
Lemma \ref{lemma-descending-property-isomorphism}.
This proves that $X$ is representable by an affine scheme and we win.
\end{proof}

\begin{lemma}
\label{lemma-descending-property-closed-immersion}
The property $\mathcal{P}(f) =$``$f$ is a closed immersion''
is fpqc local on the base.
\end{lemma}

\begin{proof}
We will use
Lemma \ref{lemma-descending-properties-morphisms}
to prove this. Assumptions (1) and (2) of that lemma follow from
Morphisms of Spaces,
Lemma \ref{spaces-morphisms-lemma-closed-immersion-local}.
Consider a cartesian diagram
$$
\xymatrix{
X' \ar[r] \ar[d] & X \ar[d] \\
Z' \ar[r] & Z
}
$$
of algebraic spaces over $S$
where $Z' \to Z$ is a surjective flat morphism of affine schemes,
and $X' \to Z'$ is a closed immersion. We have to show that $X \to Z$
is a closed immersion. The morphism $X' \to Z'$ is affine. Hence by
Lemma \ref{lemma-descending-property-affine}
we see that $X$ is a scheme and $X \to Z$ is affine.
It follows from
Descent, Lemma \ref{descent-lemma-descending-property-closed-immersion}
that $X \to Z$ is a closed immersion as desired.
\end{proof}

\begin{lemma}
\label{lemma-descending-property-separated}
The property $\mathcal{P}(f) =$``$f$ is separated''
is fpqc local on the base.
\end{lemma}

\begin{proof}
A base change of a separated morphism is separated, see
Morphisms of Spaces,
Lemma \ref{spaces-morphisms-lemma-base-change-separated}.
Hence the direct implication in
Definition \ref{definition-property-morphisms-local}.

\medskip\noindent
Let $\{Y_i \to Y\}_{i \in I}$ be an fpqc covering of algebraic spaces over $S$.
Let $f : X \to Y$ be a morphism of algebraic spaces over $S$.
Assume each base change $X_i := Y_i \times_Y X \to Y_i$ is separated.
This means that each of the morphisms
$$
\Delta_i :
X_i
\longrightarrow
X_i \times_{Y_i} X_i = Y_i \times_Y (X \times_Y X)
$$
is a closed immersion. The base change of a fpqc covering is an
fpqc covering, see
Topologies on Spaces, Lemma \ref{spaces-topologies-lemma-fpqc}
hence $\{Y_i \times_Y (X \times_Y X) \to X \times_Y X\}$
is an fpqc covering of algebraic spaces. Moreover, each
$\Delta_i$ is the base change of the morphism
$\Delta : X \to X \times_Y X$. Hence it follows from
Lemma \ref{lemma-descending-property-closed-immersion}
that $\Delta$ is a closed immersion, i.e., $f$ is separated.
\end{proof}

\begin{lemma}
\label{lemma-descending-property-proper}
The property $\mathcal{P}(f) =$``$f$ is proper''
is fpqc local on the base.
\end{lemma}

\begin{proof}
The lemma follows by combining
Lemmas \ref{lemma-descending-property-universally-closed},
\ref{lemma-descending-property-separated}
and \ref{lemma-descending-property-finite-type}.
\end{proof}

\begin{lemma}
\label{lemma-descending-property-quasi-affine}
The property $\mathcal{P}(f) =$``$f$ is quasi-affine''
is fpqc local on the base.
\end{lemma}

\begin{proof}
We will use
Lemma \ref{lemma-descending-properties-morphisms}
to prove this. Assumptions (1) and (2) of that lemma follow from
Morphisms of Spaces,
Lemma \ref{spaces-morphisms-lemma-quasi-affine-local}.
Let $Z' \to Z$ be a surjective flat morphism of affine schemes over $S$.
Let $f : X \to Z$ be a morphism of algebraic spaces, and assume
that the base change $f' : Z' \times_Z X \to Z'$ is quasi-affine.
Let $X'$ be a scheme representing $Z' \times_Z X$.
We obtain a canonical isomorphism
$$
\varphi : X' \times_Z Z' \longrightarrow Z' \times_Z X'
$$
since both schemes represent the algebraic space $Z' \times_Z Z' \times_Z X$.
This is a descent datum for $X'/Z'/Z$, see
Descent, Definition \ref{descent-definition-descent-datum}
(verification omitted, compare with
Descent, Lemma \ref{descent-lemma-descent-data-sheaves}).
Since $X' \to Z'$ is quasi-affine this descent datum is effective, see
Descent, Lemma \ref{descent-lemma-quasi-affine}.
Thus there exists a scheme $Y \to Z$ over $Z$ and an
isomorphism $\psi : Z' \times_Z Y \to X'$ compatible with descent data.
Of course $Y \to Z$ is quasi-affine (by construction or by
Descent, Lemma \ref{descent-lemma-descending-property-quasi-affine}).
Note that $\mathcal{Y} = \{Z' \times_Z Y \to Y\}$ is a
fpqc covering, and interpreting $\psi$ as an element of
$X(Z' \times_Z Y)$ we see that $\psi \in \check{H}^0(\mathcal{Y}, X)$.
By the sheaf condition for $X$ (see
Properties of Spaces, Proposition
\ref{spaces-properties-proposition-sheaf-fpqc})
we obtain a morphism $Y \to X$.
By construction the base change of this to $Z'$ is an isomorphism, hence
an isomorphism by
Lemma \ref{lemma-descending-property-isomorphism}.
This proves that $X$ is representable by a quasi-affine scheme and we win.
\end{proof}

\begin{lemma}
\label{lemma-descending-property-quasi-compact-immersion}
The property $\mathcal{P}(f) =$``$f$ is a quasi-compact immersion''
is fpqc local on the base.
\end{lemma}

\begin{proof}
We will use
Lemma \ref{lemma-descending-properties-morphisms}
to prove this. Assumptions (1) and (2) of that lemma follow from
Morphisms of Spaces,
Lemmas \ref{spaces-morphisms-lemma-closed-immersion-local} and
\ref{spaces-morphisms-lemma-quasi-compact-local}.
Consider a cartesian diagram
$$
\xymatrix{
X' \ar[r] \ar[d] & X \ar[d] \\
Z' \ar[r] & Z
}
$$
of algebraic spaces over $S$
where $Z' \to Z$ is a surjective flat morphism of affine schemes,
and $X' \to Z'$ is a quasi-compact immersion. We have to show that $X \to Z$
is a closed immersion. The morphism $X' \to Z'$ is quasi-affine. Hence by
Lemma \ref{lemma-descending-property-quasi-affine}
we see that $X$ is a scheme and $X \to Z$ is quasi-affine.
It follows from
Descent, Lemma \ref{descent-lemma-descending-property-quasi-compact-immersion}
that $X \to Z$ is a quasi-compact immersion as desired.
\end{proof}

\begin{lemma}
\label{lemma-descending-property-integral}
The property $\mathcal{P}(f) =$``$f$ is integral''
is fpqc local on the base.
\end{lemma}

\begin{proof}
An integral morphism is the same thing as an affine,
universally closed morphism. See
Morphisms of Spaces,
Lemma \ref{spaces-morphisms-lemma-integral-universally-closed}.
Hence the lemma follows on combining
Lemmas \ref{lemma-descending-property-universally-closed}
and \ref{lemma-descending-property-affine}.
\end{proof}

\begin{lemma}
\label{lemma-descending-property-finite}
The property $\mathcal{P}(f) =$``$f$ is finite''
is fpqc local on the base.
\end{lemma}

\begin{proof}
An finite morphism is the same thing as an integral,
morphism which is locally of finite type. See
Morphisms of Spaces, Lemma \ref{spaces-morphisms-lemma-finite-integral}.
Hence the lemma follows on combining
Lemmas \ref{lemma-descending-property-locally-finite-type}
and \ref{lemma-descending-property-integral}.
\end{proof}

\begin{lemma}
\label{lemma-descending-property-quasi-finite}
The properties
$\mathcal{P}(f) =$``$f$ is locally quasi-finite''
and
$\mathcal{P}(f) =$``$f$ is quasi-finite''
are fpqc local on the base.
\end{lemma}

\begin{proof}
We have already seen that ``quasi-compact'' is fpqc local on the base, see
Lemma \ref{lemma-descending-property-quasi-compact}. Hence it is enough
to prove the lemma for ``locally quasi-finite''. We will use
Lemma \ref{lemma-descending-properties-morphisms}
to prove this. Assumptions (1) and (2) of that lemma follow from
Morphisms of Spaces,
Lemma \ref{spaces-morphisms-lemma-quasi-finite-local}.
Let $Z' \to Z$ be a surjective flat morphism of affine schemes over $S$.
Let $f : X \to Z$ be a morphism of algebraic spaces, and assume
that the base change $f' : Z' \times_Z X \to Z'$ is locally quasi-finite.
We have to show that $f$ is locally quasi-finite. Let $U$ be a scheme
and let $U \to X$ be surjective and \'etale. By
Morphisms of Spaces,
Lemma \ref{spaces-morphisms-lemma-quasi-finite-local}
again, it is enough to show that $U \to Z$ is locally quasi-finite.
Since $f'$ is locally quasi-finite, and since $Z' \times_Z U$ is a
scheme \'etale over $Z' \times_Z X$ we conclude (by the same lemma again) that
$Z' \times_Z U \to Z'$ is locally quasi-finite.
As $\{Z' \to Z\}$ is an fpqc covering we conclude that
$U \to Z$ is locally quasi-finite by
Descent, Lemma \ref{descent-lemma-descending-property-quasi-finite}
as desired.
\end{proof}

\begin{lemma}
\label{lemma-descending-property-syntomic}
The property $\mathcal{P}(f) =$``$f$ is syntomic''
is fpqc local on the base.
\end{lemma}

\begin{proof}
We will use
Lemma \ref{lemma-descending-properties-morphisms}
to prove this. Assumptions (1) and (2) of that lemma follow from
Morphisms of Spaces,
Lemma \ref{spaces-morphisms-lemma-syntomic-local}.
Let $Z' \to Z$ be a surjective flat morphism of affine schemes over $S$.
Let $f : X \to Z$ be a morphism of algebraic spaces, and assume
that the base change $f' : Z' \times_Z X \to Z'$ is syntomic.
We have to show that $f$ is syntomic. Let $U$ be a scheme
and let $U \to X$ be surjective and \'etale. By
Morphisms of Spaces,
Lemma \ref{spaces-morphisms-lemma-syntomic-local}
again, it is enough to show that $U \to Z$ is syntomic.
Since $f'$ is syntomic, and since $Z' \times_Z U$ is a
scheme \'etale over $Z' \times_Z X$ we conclude (by the same lemma again) that
$Z' \times_Z U \to Z'$ is syntomic.
As $\{Z' \to Z\}$ is an fpqc covering we conclude that
$U \to Z$ is syntomic by
Descent, Lemma \ref{descent-lemma-descending-property-syntomic}
as desired.
\end{proof}

\begin{lemma}
\label{lemma-descending-property-smooth}
The property $\mathcal{P}(f) =$``$f$ is smooth''
is fpqc local on the base.
\end{lemma}

\begin{proof}
We will use
Lemma \ref{lemma-descending-properties-morphisms}
to prove this. Assumptions (1) and (2) of that lemma follow from
Morphisms of Spaces,
Lemma \ref{spaces-morphisms-lemma-smooth-local}.
Let $Z' \to Z$ be a surjective flat morphism of affine schemes over $S$.
Let $f : X \to Z$ be a morphism of algebraic spaces, and assume
that the base change $f' : Z' \times_Z X \to Z'$ is smooth.
We have to show that $f$ is smooth. Let $U$ be a scheme
and let $U \to X$ be surjective and \'etale. By
Morphisms of Spaces,
Lemma \ref{spaces-morphisms-lemma-smooth-local}
again, it is enough to show that $U \to Z$ is smooth.
Since $f'$ is smooth, and since $Z' \times_Z U$ is a
scheme \'etale over $Z' \times_Z X$ we conclude (by the same lemma again) that
$Z' \times_Z U \to Z'$ is smooth.
As $\{Z' \to Z\}$ is an fpqc covering we conclude that
$U \to Z$ is smooth by
Descent, Lemma \ref{descent-lemma-descending-property-smooth}
as desired.
\end{proof}

\begin{lemma}
\label{lemma-descending-property-unramified}
The property $\mathcal{P}(f) =$``$f$ is unramified''
is fpqc local on the base.
\end{lemma}

\begin{proof}
We will use
Lemma \ref{lemma-descending-properties-morphisms}
to prove this. Assumptions (1) and (2) of that lemma follow from
Morphisms of Spaces,
Lemma \ref{spaces-morphisms-lemma-unramified-local}.
Let $Z' \to Z$ be a surjective flat morphism of affine schemes over $S$.
Let $f : X \to Z$ be a morphism of algebraic spaces, and assume
that the base change $f' : Z' \times_Z X \to Z'$ is unramified.
We have to show that $f$ is unramified. Let $U$ be a scheme
and let $U \to X$ be surjective and \'etale. By
Morphisms of Spaces,
Lemma \ref{spaces-morphisms-lemma-unramified-local}
again, it is enough to show that $U \to Z$ is unramified.
Since $f'$ is unramified, and since $Z' \times_Z U$ is a
scheme \'etale over $Z' \times_Z X$ we conclude (by the same lemma again) that
$Z' \times_Z U \to Z'$ is unramified.
As $\{Z' \to Z\}$ is an fpqc covering we conclude that
$U \to Z$ is unramified by
Descent, Lemma \ref{descent-lemma-descending-property-unramified}
as desired.
\end{proof}

\begin{lemma}
\label{lemma-descending-property-etale}
The property $\mathcal{P}(f) =$``$f$ is \'etale''
is fpqc local on the base.
\end{lemma}

\begin{proof}
We will use
Lemma \ref{lemma-descending-properties-morphisms}
to prove this. Assumptions (1) and (2) of that lemma follow from
Morphisms of Spaces,
Lemma \ref{spaces-morphisms-lemma-etale-local}.
Let $Z' \to Z$ be a surjective flat morphism of affine schemes over $S$.
Let $f : X \to Z$ be a morphism of algebraic spaces, and assume
that the base change $f' : Z' \times_Z X \to Z'$ is \'etale.
We have to show that $f$ is \'etale. Let $U$ be a scheme
and let $U \to X$ be surjective and \'etale. By
Morphisms of Spaces,
Lemma \ref{spaces-morphisms-lemma-etale-local}
again, it is enough to show that $U \to Z$ is \'etale.
Since $f'$ is \'etale, and since $Z' \times_Z U$ is a
scheme \'etale over $Z' \times_Z X$ we conclude (by the same lemma again) that
$Z' \times_Z U \to Z'$ is \'etale.
As $\{Z' \to Z\}$ is an fpqc covering we conclude that
$U \to Z$ is \'etale by
Descent, Lemma \ref{descent-lemma-descending-property-etale}
as desired.
\end{proof}

\begin{lemma}
\label{lemma-descending-property-finite-locally-free}
The property $\mathcal{P}(f) =$``$f$ is finite locally free''
is fpqc local on the base.
\end{lemma}

\begin{proof}
Being finite locally free is equivalent to being
finite, flat and locally of finite presentation
(Morphisms of Spaces, Lemma \ref{spaces-morphisms-lemma-finite-flat}).
Hence this follows from Lemmas
\ref{lemma-descending-property-finite},
\ref{lemma-descending-property-flat}, and
\ref{lemma-descending-property-locally-finite-presentation}.
\end{proof}

\begin{lemma}
\label{lemma-descending-property-monomorphism}
The property $\mathcal{P}(f) =$``$f$ is a monomorphism''
is fpqc local on the base.
\end{lemma}

\begin{proof}
Let $f : X \to Y$ be a morphism of algebraic spaces.
Let $\{Y_i \to Y\}$ be an fpqc covering, and assume
each of the base changes $f_i : X_i \to Y_i$ of $f$ is
a monomorphism. We have to show that $f$ is a monomorphism.

\medskip\noindent
First proof. Note that $f$ is a monomorphism if and only if
$\Delta : X \to X \times_Y X$ is an isomorphism. By applying this to
$f_i$ we see that each of the morphisms
$$
\Delta_i :
X_i
\longrightarrow
X_i \times_{Y_i} X_i = Y_i \times_Y (X \times_Y X)
$$
is an isomorphism. The base change of an fpqc covering is an fpqc covering, see
Topologies on Spaces, Lemma \ref{spaces-topologies-lemma-fpqc}
hence $\{Y_i \times_Y (X \times_Y X) \to X \times_Y X\}$
is an fpqc covering of algebraic spaces. Moreover, each
$\Delta_i$ is the base change of the morphism
$\Delta : X \to X \times_Y X$. Hence it follows from
Lemma \ref{lemma-descending-property-isomorphism}
that $\Delta$ is an isomorphism, i.e., $f$ is a monomorphism.

\medskip\noindent
Second proof.
Let $V$ be a scheme, and let $V \to Y$ be a surjective \'etale morphism.
If we can show that $V \times_Y X \to V$ is a monomorphism, then it
follows that $X \to Y$ is a monomorphism. Namely, given any
cartesian diagram of sheaves
$$
\vcenter{
\xymatrix{
\mathcal{F} \ar[r]_a \ar[d]_b & \mathcal{G} \ar[d]^c \\
\mathcal{H} \ar[r]^d & \mathcal{I}
}
}
\quad
\quad
\mathcal{F} = \mathcal{H} \times_\mathcal{I} \mathcal{G}
$$
if $c$ is a surjection of sheaves, and $a$ is injective, then also
$d$ is injective. This reduces the problem to the case where $Y$ is
a scheme. Moreover, in this case we may assume that the algebraic spaces
$Y_i$ are schemes also, since we can always refine the covering to place
ourselves in this situation, see
Topologies on Spaces, Lemma \ref{spaces-topologies-lemma-refine-fpqc-schemes}.

\medskip\noindent
Assume $\{Y_i \to Y\}$ is an fpqc covering of schemes.
Let $a, b : T \to X$ be two morphisms
such that $f \circ a = f \circ b$. We have to show that $a = b$.
Since $f_i$ is a monomorphism we see that $a_i = b_i$, where
$a_i, b_i : Y_i \times_Y T \to X_i$ are
the base changes. In particular the compositions
$Y_i \times_Y T \to T \to X$ are equal.
Since $\{Y_i \times_Y T \to T\}$ is an fpqc covering we
deduce that $a = b$ from Properties of Spaces, Proposition
\ref{spaces-properties-proposition-sheaf-fpqc}.
\end{proof}




\section{Descending properties of morphisms in the fppf topology}
\label{section-descending-properties-morphisms-fppf}

\noindent
In this section we find some properties of morphisms of algebraic spaces
for which we could not (yet) show they are local on the base in
the fpqc topology which, however, are local on the base
in the fppf topology.

\begin{lemma}
\label{lemma-descending-fppf-property-immersion}
The property $\mathcal{P}(f) =$``$f$ is an immersion''
is fppf local on the base.
\end{lemma}

\begin{proof}
Let $f : X \to Y$ be a morphism of algebraic spaces.
Let $\{Y_i \to Y\}_{i \in I}$ be an fppf covering of $Y$.
Let $f_i : X_i \to Y_i$ be the base change of $f$.

\medskip\noindent
If $f$ is an immersion, then each $f_i$ is an immersion by
Spaces, Lemma \ref{spaces-lemma-base-change-immersions}.
This proves the direct implication in
Definition \ref{definition-property-morphisms-local}.

\medskip\noindent
Conversely, assume each $f_i$ is an immersion. By
Morphisms of Spaces,
Lemma \ref{spaces-morphisms-lemma-immersions-monomorphisms}
this implies each $f_i$ is separated. By
Morphisms of Spaces,
Lemma \ref{spaces-morphisms-lemma-immersion-quasi-finite}
this implies each $f_i$ is locally quasi-finite.
Hence we see that $f$ is locally quasi-finite and separated, by applying
Lemmas \ref{lemma-descending-property-separated}
and \ref{lemma-descending-property-quasi-finite}.
By
Morphisms of Spaces, Lemma
\ref{spaces-morphisms-lemma-locally-quasi-finite-separated-representable}
this implies that $f$ is representable!

\medskip\noindent
By
Morphisms of Spaces, Lemma \ref{spaces-morphisms-lemma-closed-immersion-local}
it suffices to show that for every scheme $Z$ and morphism $Z \to Y$
the base change $Z \times_Y X \to Z$ is an immersion. By
Topologies on Spaces, Lemma \ref{spaces-topologies-lemma-refine-fppf-schemes}
we can find an fppf covering $\{Z_i \to Z\}$ by schemes which refines
the pullback of the covering $\{Y_i \to Y\}$ to $Z$.
Hence we see that $Z \times_Y X \to Z$ (which is a morphism of schemes
according to the result of the preceding paragraph) becomes an immersion
after pulling back to the members of an fppf (by schemes) of $Z$.
Hence $Z \times_Y X \to Z$ is an immersion by the result for schemes, see
Descent, Lemma \ref{descent-lemma-descending-fppf-property-immersion}.
\end{proof}

\begin{lemma}
\label{lemma-descending-fppf-property-locally-separated}
The property $\mathcal{P}(f) =$``$f$ is locally separated''
is fppf local on the base.
\end{lemma}

\begin{proof}
A base change of a locally separated morphism is locally separated, see
Morphisms of Spaces,
Lemma \ref{spaces-morphisms-lemma-base-change-separated}.
Hence the direct implication in
Definition \ref{definition-property-morphisms-local}.

\medskip\noindent
Let $\{Y_i \to Y\}_{i \in I}$ be an fppf covering of algebraic spaces over $S$.
Let $f : X \to Y$ be a morphism of algebraic spaces over $S$.
Assume each base change $X_i := Y_i \times_Y X \to Y_i$ is locally separated.
This means that each of the morphisms
$$
\Delta_i :
X_i
\longrightarrow
X_i \times_{Y_i} X_i = Y_i \times_Y (X \times_Y X)
$$
is an immersion. The base change of a fppf covering is an
fppf covering, see
Topologies on Spaces, Lemma \ref{spaces-topologies-lemma-fppf}
hence $\{Y_i \times_Y (X \times_Y X) \to X \times_Y X\}$
is an fppf covering of algebraic spaces. Moreover, each
$\Delta_i$ is the base change of the morphism
$\Delta : X \to X \times_Y X$. Hence it follows from
Lemma \ref{lemma-descending-fppf-property-immersion}
that $\Delta$ is an immersion, i.e., $f$ is locally separated.
\end{proof}





\section{Application of descent of properties of morphisms}
\label{section-application-descending-properties-morphisms}

\noindent
This section is the analogue of Descent, Section
\ref{descent-section-application-descending-properties-morphisms}.

\begin{lemma}
\label{lemma-descending-property-ample}
Let $S$ be a scheme.
Let $f : X \to Y$ be a morphism of algebraic spaces over $S$.
Let $\mathcal{L}$ be an invertible $\mathcal{O}_X$-module.
Let $\{g_i : Y_i \to Y\}_{i \in I}$ be an fpqc covering.
Let $f_i : X_i \to Y_i$ be the base change of $f$ and let $\mathcal{L}_i$
be the pullback of $\mathcal{L}$ to $X_i$.
The following are equivalent
\begin{enumerate}
\item $\mathcal{L}$ is ample on $X/Y$, and
\item $\mathcal{L}_i$ is ample on $X_i/Y_i$
for every $i \in I$.
\end{enumerate}
\end{lemma}

\begin{proof}
The implication (1) $\Rightarrow$ (2) follows from
Divisors on Spaces, Lemma \ref{spaces-divisors-lemma-ample-base-change}.
Assume (2). To check $\mathcal{L}$ is ample on $X/Y$ we may
work \'etale locally on $Y$, see
Divisors on Spaces, Lemma \ref{spaces-divisors-lemma-relatively-ample-local}.
Thus we may assume that $Y$ is a scheme and then we may
in turn assume each $Y_i$ is a scheme too, see
Topologies on Spaces, Lemma \ref{spaces-topologies-lemma-refine-fpqc-schemes}.
In other words, we may assume that
$\{Y_i \to Y\}$ is an fpqc covering of schemes.

\medskip\noindent
By Divisors on Spaces, Lemma
\ref{spaces-divisors-lemma-relatively-ample-properties}
we see that $X_i \to Y_i$ is representable (i.e., $X_i$ is a scheme),
quasi-compact, and separated. Hence $f$ is quasi-compact and separated by
Lemmas \ref{lemma-descending-property-quasi-compact} and
\ref{lemma-descending-property-separated}.
This means that
$\mathcal{A} = \bigoplus_{d \geq 0} f_*\mathcal{L}^{\otimes d}$
is a quasi-coherent graded $\mathcal{O}_Y$-algebra
(Morphisms of Spaces, Lemma \ref{spaces-morphisms-lemma-pushforward}).
Moreover, the formation of $\mathcal{A}$ commutes with flat
base change by
Cohomology of Spaces, Lemma
\ref{spaces-cohomology-lemma-flat-base-change-cohomology}.
In particular, if we set
$\mathcal{A}_i = \bigoplus_{d \geq 0} f_{i, *}\mathcal{L}_i^{\otimes d}$
then we have $\mathcal{A}_i = g_i^*\mathcal{A}$.
It follows that the natural maps
$\psi_d : f^*\mathcal{A}_d \to \mathcal{L}^{\otimes d}$
of $\mathcal{O}_X$
pullback to give the natural maps
$\psi_{i, d} : f_i^*(\mathcal{A}_i)_d \to \mathcal{L}_i^{\otimes d}$
of $\mathcal{O}_{X_i}$-modules. Since $\mathcal{L}_i$ is ample on $X_i/Y_i$
we see that for any point $x_i \in X_i$, there exists a $d \geq 1$
such that $f_i^*(\mathcal{A}_i)_d \to \mathcal{L}_i^{\otimes d}$
is surjective on stalks at $x_i$. This follows either directly
from the definition of a relatively ample module or from
Morphisms, Lemma \ref{morphisms-lemma-characterize-relatively-ample}.
If $x \in |X|$, then we can choose an $i$ and an $x_i \in X_i$
mapping to $x$. Since
$\mathcal{O}_{X, \overline{x}} \to \mathcal{O}_{X_i, \overline{x}_i}$
is flat hence faithfully flat, we conclude that for every $x \in |X|$
there exists a $d \geq 1$ such that
$f^*\mathcal{A}_d \to \mathcal{L}^{\otimes d}$
is surjective on stalks at $x$.
This implies that the open subset $U(\psi) \subset X$ of
Divisors on Spaces, Lemma
\ref{spaces-divisors-lemma-invertible-map-into-relative-proj}
corresponding to the map
$\psi : f^*\mathcal{A} \to \bigoplus_{d \geq 0} \mathcal{L}^{\otimes d}$
of graded $\mathcal{O}_X$-algebras
is equal to $X$. Consider the corresponding morphism
$$
r_{\mathcal{L}, \psi} : X \longrightarrow \underline{\text{Proj}}_Y(\mathcal{A})
$$
It is clear from the above that the base change of
$r_{\mathcal{L}, \psi}$ to $Y_i$ is the morphism
$r_{\mathcal{L}_i, \psi_i}$ which is an open immersion by
Morphisms, Lemma \ref{morphisms-lemma-characterize-relatively-ample}.
Hence $r_{\mathcal{L}, \psi}$ is an open immersion
by Lemma \ref{lemma-descending-property-open-immersion}.
Hence $X$ is a scheme
and we conclude $\mathcal{L}$ is ample on $X/Y$ by
Morphisms, Lemma \ref{morphisms-lemma-characterize-relatively-ample}.
\end{proof}

\begin{lemma}
\label{lemma-ample-in-neighbourhood}
Let $S$ be a scheme.
Let $f : X \to Y$ be a proper morphism of algebraic spaces over $S$.
Let $\mathcal{L}$ be an invertible $\mathcal{O}_X$-module.
There exists an open subspace $V \subset Y$ characterized by
the following property:
A morphism $Y' \to Y$ of algebraic spaces factors
through $V$ if and only if the pullback $\mathcal{L}'$
of $\mathcal{L}$ to $X' = Y' \times_Y X$ is ample on $X'/Y'$
(as in Divisors on Spaces, Definition
\ref{spaces-divisors-definition-relatively-ample}).
\end{lemma}

\begin{proof}
Suppose that the lemma holds whenever $Y$ is a scheme.
Let $U$ be a scheme and let $U \to Y$ be a surjective \'etale morphism.
Let $R = U \times_Y U$ with projections $t, s : R \to U$.
Denote $X_U = U \times_Y X$ and $\mathcal{L}_U$ the pullback.
Then we get an open subscheme $V' \subset U$ as in the lemma
for $(X_U \to U, \mathcal{L}_U)$. By the functorial characterization
we see that $s^{-1}(V') = t^{-1}(V')$.
Thus there is an open subspace $V \subset Y$ such that
$V'$ is the inverse image of $V$ in $U$.
In particular $V' \to V$ is surjective \'etale and we
conclude that $\mathcal{L}_V$ is ample on $X_V/V$
(Divisors on Spaces, Lemma \ref{spaces-divisors-lemma-relatively-ample-local}).
Now, if $Y' \to Y$ is a morphism such that
$\mathcal{L}'$ is ample on $X'/Y'$, then
$U \times_Y Y' \to Y'$ must factor through $V'$
and we conclude that $Y' \to Y$ factors through $V$.
Hence $V \subset Y$ is as in the statement of the lemma.
In this way we reduce to the case dealt with in the next
paragraph.

\medskip\noindent
Assume $Y$ is a scheme. Since the question is local on $Y$
we may assume $Y$ is an affine scheme. We will show the
following:
\begin{enumerate}
\item[(A)] If $\Spec(k) \to Y$ is a morphism such that
$\mathcal{L}_k$ is ample on $X_k/k$, then there is an
open neighbourhood $V \subset Y$ of the image of $\Spec(k) \to Y$
such that $\mathcal{L}_V$ is ample on $X_V/V$.
\end{enumerate}
It is clear that (A) implies the truth of the lemma.

\medskip\noindent
Let $X \to Y$, $\mathcal{L}$, $\Spec(k) \to Y$ be as in (A).
By Lemma \ref{lemma-descending-property-ample}
we may assume that $k = \kappa(y)$ is the residue field of a point $y$
of $Y$.

\medskip\noindent
As $Y$ is affine we can find a directed set $I$ and an
inverse system of morphisms $X_i \to Y_i$ of algebraic spaces
with $Y_i$ of finite presentation over $\mathbf{Z}$, with affine
transition morphisms $X_i \to X_{i'}$ and $Y_i \to Y_{i'}$,
with $X_i \to Y_i$ proper and of finite presentation, and such that
$X \to Y = \lim (X_i \to Y_i)$. See Limits of Spaces, Lemma
\ref{spaces-limits-lemma-proper-limit-of-proper-finite-presentation-noetherian}.
After shrinking $I$ we may assume $Y_i$ is an (affine) scheme for all $i$,
see Limits of Spaces, Lemma \ref{spaces-limits-lemma-limit-is-affine}.
After shrinking $I$ we can assume we have a compatible system of
invertible $\mathcal{O}_{X_i}$-modules $\mathcal{L}_i$
pulling back to $\mathcal{L}$, see
Limits of Spaces, Lemma \ref{spaces-limits-lemma-descend-invertible-modules}.
Let $y_i \in Y_i$ be the image of $y$.
Then $\kappa(y) = \colim \kappa(y_i)$.
Hence $X_y = \lim X_{i, y_i}$ and after shrinking $I$ we
may assume $X_{i, y_i}$ is a scheme for all $i$, see
Limits of Spaces, Lemma \ref{spaces-limits-lemma-limit-is-scheme}.
Hence for some $i$ we have $\mathcal{L}_{i, y_i}$
is ample on $X_{i, y_i}$ by
Limits, Lemma \ref{limits-lemma-limit-ample}.
By Divisors on Spaces, Lemma \ref{spaces-divisors-lemma-ample-in-neighbourhood}
we find an open neighbourhood
$V_i \subset Y_i$ of $y_i$ such that
$\mathcal{L}_i$ restricted to $f_i^{-1}(V_i)$
is ample relative to $V_i$.
Letting $V \subset Y$ be the inverse image of
$V_i$ finishes the proof (hints: use
Morphisms, Lemma \ref{morphisms-lemma-ample-base-change} and
the fact that $X \to Y \times_{Y_i} X_i$ is affine
and the fact that the pullback of an
ample invertible sheaf by an affine morphism is ample by
Morphisms, Lemma \ref{morphisms-lemma-pullback-ample-tensor-relatively-ample}).
\end{proof}













\section{Properties of morphisms local on the source}
\label{section-properties-morphisms-local-source}

\noindent
In this section we define what it means for a property of morphisms of
algebraic spaces to be local on the source. Please compare with
Descent, Section \ref{descent-section-properties-morphisms-local-source}.

\begin{definition}
\label{definition-property-morphisms-local-source}
Let $S$ be a scheme.
Let $\mathcal{P}$ be a property of morphisms of algebraic spaces over $S$.
Let $\tau \in \{fpqc, \linebreak[0] fppf, \linebreak[0] syntomic, \linebreak[0]
smooth, \linebreak[0] \etale\}$. We say $\mathcal{P}$ is
{\it $\tau$ local on the source}, or
{\it local on the source for the $\tau$-topology} if for
any morphism $f : X \to Y$ of algebraic spaces over $S$, and any
$\tau$-covering $\{X_i \to X\}_{i \in I}$ of algebraic spaces we have
$$
f \text{ has }\mathcal{P}
\Leftrightarrow
\text{each }X_i \to Y\text{ has }\mathcal{P}.
$$
\end{definition}

\noindent
To be sure, since isomorphisms are always coverings
we see (or require) that property $\mathcal{P}$ holds for $X \to Y$
if and only if it holds for any arrow $X' \to Y'$ isomorphic to $X \to Y$.
If a property is $\tau$-local on the source then it is preserved by
precomposing with morphisms which occur in $\tau$-coverings. Here
is a formal statement.

\begin{lemma}
\label{lemma-precompose-property-local-source}
Let $S$ be a scheme.
Let $\tau \in \{fpqc, \linebreak[0] fppf, \linebreak[0] syntomic, \linebreak[0]
smooth, \linebreak[0] \etale\}$.
Let $\mathcal{P}$ be a property of morphisms of algebraic spaces over $S$
which is $\tau$ local on the source. Let $f : X \to Y$ have property
$\mathcal{P}$. For any morphism $a : X' \to X$ which is
flat, resp.\ flat and locally of finite presentation, resp.\ syntomic,
resp.\ smooth, resp.\ \'etale, the composition $f \circ a : X' \to Y$ has
property $\mathcal{P}$.
\end{lemma}

\begin{proof}
This is true because we can fit $X' \to X$ into a family of
morphisms which forms a $\tau$-covering.
\end{proof}

\begin{lemma}
\label{lemma-transfer-from-schemes}
Let $S$ be a scheme.
Let $\tau \in \{fpqc, \linebreak[0] fppf, \linebreak[0] syntomic, \linebreak[0]
smooth, \linebreak[0] \etale\}$.
Suppose that $\mathcal{P}$ is a property of morphisms of schemes over $S$
which is \'etale local on the source-and-target. Denote $\mathcal{P}_{spaces}$
the corresponding property of morphisms of algebraic spaces over $S$, see
Morphisms of Spaces, Definition \ref{spaces-morphisms-definition-P}.
If $\mathcal{P}$ is local on the source for the $\tau$-topology, then
$\mathcal{P}_{spaces}$ is local on the source for the $\tau$-topology.
\end{lemma}

\begin{proof}
Let $f : X \to Y$ be a morphism of algebraic spaces over $S$.
Let $\{X_i \to X\}_{i \in I}$ be a $\tau$-covering of algebraic spaces.
Choose a scheme $V$ and a surjective \'etale morphism $V \to Y$.
Choose a scheme $U$ and a surjective \'etale morphism $U \to X \times_Y V$.
For each $i$ choose a scheme $U_i$ and a surjective \'etale morphism
$U_i \to X_i \times_X U$.

\medskip\noindent
Note that $\{X_i \times_X U \to U\}_{i \in I}$ is a $\tau$-covering.
Note that each $\{U_i \to X_i \times_X U\}$ is an \'etale covering,
hence a $\tau$-covering. Hence $\{U_i \to U\}_{i \in I}$ is a
$\tau$-covering of algebraic spaces over $S$. But since $U$ and each $U_i$
is a scheme we see that $\{U_i \to U\}_{i \in I}$ is a
$\tau$-covering of schemes over $S$.

\medskip\noindent
Now we have
\begin{align*}
f \text{ has }\mathcal{P}_{spaces}
& \Leftrightarrow
U \to V \text{ has }\mathcal{P} \\
& \Leftrightarrow
\text{each }U_i \to V \text{ has }\mathcal{P} \\
& \Leftrightarrow
\text{each }X_i \to Y\text{ has }\mathcal{P}_{spaces}.
\end{align*}
the first and last equivalence by the definition of
$\mathcal{P}_{spaces}$ the middle equivalence because we assumed
$\mathcal{P}$ is local on the source in the $\tau$-topology.
\end{proof}






\section{Properties of morphisms local in the fpqc topology on the source}
\label{section-fpqc-local-source}

\noindent
Here are some properties of morphisms that are fpqc local on the source.

\begin{lemma}
\label{lemma-flat-fpqc-local-source}
The property $\mathcal{P}(f)=$``$f$ is flat'' is fpqc local on the source.
\end{lemma}

\begin{proof}
Follows from
Lemma \ref{lemma-transfer-from-schemes}
using
Morphisms of Spaces, Definition \ref{spaces-morphisms-definition-flat}
and
Descent, Lemma \ref{descent-lemma-flat-fpqc-local-source}.
\end{proof}










\section{Properties of morphisms local in the fppf topology on the source}
\label{section-fppf-local-source}

\noindent
Here are some properties of morphisms that are fppf local on the source.

\begin{lemma}
\label{lemma-locally-finite-presentation-fppf-local-source}
The property $\mathcal{P}(f)=$``$f$ is locally of finite presentation''
is fppf local on the source.
\end{lemma}

\begin{proof}
Follows from
Lemma \ref{lemma-transfer-from-schemes}
using
Morphisms of Spaces,
Definition \ref{spaces-morphisms-definition-locally-finite-presentation}
and
Descent,
Lemma \ref{descent-lemma-locally-finite-presentation-fppf-local-source}.
\end{proof}

\begin{lemma}
\label{lemma-locally-finite-type-fppf-local-source}
The property $\mathcal{P}(f)=$``$f$ is locally of finite type''
is fppf local on the source.
\end{lemma}

\begin{proof}
Follows from
Lemma \ref{lemma-transfer-from-schemes}
using
Morphisms of Spaces,
Definition \ref{spaces-morphisms-definition-locally-finite-type}
and
Descent, Lemma \ref{descent-lemma-locally-finite-type-fppf-local-source}.
\end{proof}

\begin{lemma}
\label{lemma-open-fppf-local-source}
The property $\mathcal{P}(f)=$``$f$ is open''
is fppf local on the source.
\end{lemma}

\begin{proof}
Follows from
Lemma \ref{lemma-transfer-from-schemes}
using
Morphisms of Spaces, Definition \ref{spaces-morphisms-definition-open}
and
Descent, Lemma \ref{descent-lemma-open-fppf-local-source}.
\end{proof}

\begin{lemma}
\label{lemma-universally-open-fppf-local-source}
The property $\mathcal{P}(f)=$``$f$ is universally open''
is fppf local on the source.
\end{lemma}

\begin{proof}
Follows from
Lemma \ref{lemma-transfer-from-schemes}
using
Morphisms of Spaces, Definition \ref{spaces-morphisms-definition-open}
and
Descent, Lemma \ref{descent-lemma-universally-open-fppf-local-source}.
\end{proof}



\section{Properties of morphisms local in the syntomic topology on the source}
\label{section-syntomic-local-source}

\noindent
Here are some properties of morphisms that are syntomic local on the source.

\begin{lemma}
\label{lemma-syntomic-syntomic-local-source}
The property $\mathcal{P}(f)=$``$f$ is syntomic''
is syntomic local on the source.
\end{lemma}

\begin{proof}
Follows from
Lemma \ref{lemma-transfer-from-schemes}
using
Morphisms of Spaces, Definition \ref{spaces-morphisms-definition-syntomic}
and
Descent, Lemma \ref{descent-lemma-syntomic-syntomic-local-source}.
\end{proof}




\section{Properties of morphisms local in the smooth topology on the source}
\label{section-smooth-local-source}

\noindent
Here are some properties of morphisms that are smooth local on the source.

\begin{lemma}
\label{lemma-smooth-smooth-local-source}
The property $\mathcal{P}(f)=$``$f$ is smooth''
is smooth local on the source.
\end{lemma}

\begin{proof}
Follows from
Lemma \ref{lemma-transfer-from-schemes}
using
Morphisms of Spaces, Definition \ref{spaces-morphisms-definition-smooth}
and
Descent, Lemma \ref{descent-lemma-smooth-smooth-local-source}.
\end{proof}



\section{Properties of morphisms local in the \'etale topology on the source}
\label{section-etale-local-source}

\noindent
Here are some properties of morphisms that are \'etale local on the source.

\begin{lemma}
\label{lemma-etale-etale-local-source}
The property $\mathcal{P}(f)=$``$f$ is \'etale''
is \'etale local on the source.
\end{lemma}

\begin{proof}
Follows from
Lemma \ref{lemma-transfer-from-schemes}
using
Morphisms of Spaces,
Definition \ref{spaces-morphisms-definition-etale}
and
Descent, Lemma \ref{descent-lemma-etale-etale-local-source}.
\end{proof}

\begin{lemma}
\label{lemma-locally-quasi-finite-etale-local-source}
The property $\mathcal{P}(f)=$``$f$ is locally quasi-finite''
is \'etale local on the source.
\end{lemma}

\begin{proof}
Follows from
Lemma \ref{lemma-transfer-from-schemes}
using
Morphisms of Spaces,
Definition \ref{spaces-morphisms-definition-locally-quasi-finite}
and
Descent, Lemma \ref{descent-lemma-locally-quasi-finite-etale-local-source}.
\end{proof}

\begin{lemma}
\label{lemma-unramified-etale-local-source}
The property $\mathcal{P}(f)=$``$f$ is unramified''
is \'etale local on the source.
\end{lemma}

\begin{proof}
Follows from
Lemma \ref{lemma-transfer-from-schemes}
using
Morphisms of Spaces, Definition \ref{spaces-morphisms-definition-unramified}
and
Descent, Lemma \ref{descent-lemma-unramified-etale-local-source}.
\end{proof}




\section{Properties of morphisms smooth local on source-and-target}
\label{section-properties-local-source-target}

\noindent
Let $\mathcal{P}$ be a property of morphisms of algebraic spaces. There is an
intuitive meaning to the phrase ``$\mathcal{P}$ is smooth local on the
source and target''. However, it turns out that this notion is not
the same as asking $\mathcal{P}$ to be both smooth
local on the source and smooth local on the target.
We have discussed a similar phenomenon (for the \'etale topology and
the category of schemes) in great detail in
Descent, Section \ref{descent-section-properties-etale-local-source-target}
(for a quick overview take a look at
Descent, Remark \ref{descent-remark-compare-definitions}).
However, there is an important difference between the case of the smooth
and the \'etale topology. To see this difference we encourage the reader
to ponder the difference between
Descent, Lemma \ref{descent-lemma-local-source-target-implies}
and
Lemma \ref{lemma-local-source-target-implies}
as well as the difference between
Descent, Lemma \ref{descent-lemma-local-source-target-characterize}
and
Lemma \ref{lemma-local-source-target-characterize}.
Namely, in the \'etale setting the choice of the \'etale ``covering'' of the
target is immaterial, whereas in the smooth setting it is not.

\begin{definition}
\label{definition-local-source-target}
Let $S$ be a scheme.
Let $\mathcal{P}$ be a property of morphisms of algebraic spaces over $S$.
We say $\mathcal{P}$ is {\it smooth local on source-and-target} if
\begin{enumerate}
\item (stable under precomposing with smooth maps)
if $f : X \to Y$ is smooth and $g : Y \to Z$ has $\mathcal{P}$,
then $g \circ f$ has $\mathcal{P}$,
\item (stable under smooth base change)
if $f : X \to Y$ has $\mathcal{P}$ and $Y' \to Y$ is smooth, then
the base change $f' : Y' \times_Y X \to Y'$ has $\mathcal{P}$, and
\item (locality) given a morphism $f : X \to Y$ the following are
equivalent
\begin{enumerate}
\item $f$ has $\mathcal{P}$,
\item for every $x \in |X|$ there exists a commutative diagram
$$
\xymatrix{
U \ar[d]_a \ar[r]_h & V \ar[d]^b \\
X \ar[r]^f & Y
}
$$
with smooth vertical arrows and $u \in |U|$ with $a(u) = x$ such that
$h$ has $\mathcal{P}$.
\end{enumerate}
\end{enumerate}
\end{definition}

\noindent
The above serves as our definition. In the lemmas below we will show that
this is equivalent to $\mathcal{P}$ being smooth local on the target,
smooth local on the source, and stable under post-composing by smooth morphisms.

\begin{lemma}
\label{lemma-local-source-target-implies}
Let $S$ be a scheme.
Let $\mathcal{P}$ be a property of morphisms of algebraic spaces over $S$
which is smooth local on source-and-target. Then
\begin{enumerate}
\item $\mathcal{P}$ is smooth local on the source,
\item $\mathcal{P}$ is smooth local on the target,
\item $\mathcal{P}$ is stable under postcomposing with smooth morphisms:
if $f : X \to Y$ has $\mathcal{P}$ and $g : Y \to Z$ is smooth, then
$g \circ f$ has $\mathcal{P}$.
\end{enumerate}
\end{lemma}

\begin{proof}
We write everything out completely.

\medskip\noindent
Proof of (1). Let $f : X \to Y$ be a morphism of algebraic spaces over $S$.
Let $\{X_i \to X\}_{i \in I}$ be a smooth covering of $X$. If each composition
$h_i : X_i \to Y$ has $\mathcal{P}$, then for each $|x| \in X$ we can find
an $i \in I$ and a point $x_i \in |X_i|$ mapping to $x$. Then
$(X_i, x_i) \to (X, x)$ is a smooth morphism of pairs, and
$\text{id}_Y : Y \to Y$ is a smooth morphism, and $h_i$ is as in part (3) of
Definition \ref{definition-local-source-target}.
Thus we see that $f$ has $\mathcal{P}$.
Conversely, if $f$ has $\mathcal{P}$ then each $X_i \to Y$ has
$\mathcal{P}$ by
Definition \ref{definition-local-source-target} part (1).

\medskip\noindent
Proof of (2). Let $f : X \to Y$ be a morphism of algebraic spaces over $S$.
Let $\{Y_i \to Y\}_{i \in I}$ be a smooth covering of $Y$.
Write $X_i = Y_i \times_Y X$ and $h_i : X_i \to Y_i$ for the base change
of $f$.  If each  $h_i : X_i \to Y_i$ has $\mathcal{P}$, then for each
$x \in |X|$ we pick an $i \in I$ and a point $x_i \in |X_i|$ mapping to $x$.
Then $(X_i, x_i) \to (X, x)$ is a smooth morphism of pairs, $Y_i \to Y$ is
smooth, and $h_i$ is as in part (3) of
Definition \ref{definition-local-source-target}.
Thus we see that $f$ has $\mathcal{P}$.
Conversely, if $f$ has $\mathcal{P}$, then each $X_i \to Y_i$ has
$\mathcal{P}$ by
Definition \ref{definition-local-source-target} part (2).

\medskip\noindent
Proof of (3). Assume $f : X \to Y$ has $\mathcal{P}$ and $g : Y \to Z$ is
smooth. For every $x \in |X|$ we can think of $(X, x) \to (X, x)$ as a
smooth morphism of pairs, $Y \to Z$ is a smooth morphism, and $h = f$ is as
in part (3) of
Definition \ref{definition-local-source-target}.
Thus we see that $g \circ f$ has $\mathcal{P}$.
\end{proof}

\noindent
The following lemma is the analogue of
Morphisms, Lemma \ref{morphisms-lemma-locally-P-characterize}.

\begin{lemma}
\label{lemma-local-source-target-characterize}
Let $S$ be a scheme.
Let $\mathcal{P}$ be a property of morphisms of algebraic spaces over $S$
which is smooth local on source-and-target. Let $f : X \to Y$ be a morphism
of algebraic spaces over $S$. The following are equivalent:
\begin{enumerate}
\item[(a)] $f$ has property $\mathcal{P}$,
\item[(b)] for every $x \in |X|$ there exists a smooth morphism of pairs
$a : (U, u) \to (X, x)$, a smooth morphism $b : V \to Y$, and
a morphism $h : U \to V$ such that $f \circ a = b \circ h$ and
$h$ has $\mathcal{P}$,
\item[(c)] for some commutative diagram
$$
\xymatrix{
U \ar[d]_a \ar[r]_h & V \ar[d]^b \\
X \ar[r]^f & Y
}
$$
with $a$, $b$ smooth and $a$ surjective the morphism $h$ has $\mathcal{P}$,
\item[(d)] for any commutative diagram
$$
\xymatrix{
U \ar[d]_a \ar[r]_h & V \ar[d]^b \\
X \ar[r]^f & Y
}
$$
with $b$ smooth and $U \to X \times_Y V$ smooth
the morphism $h$ has $\mathcal{P}$,
\item[(e)] there exists a smooth covering $\{Y_i \to Y\}_{i \in I}$ such
that each base change $Y_i \times_Y X \to Y_i$ has $\mathcal{P}$,
\item[(f)] there exists a smooth covering $\{X_i \to X\}_{i \in I}$ such
that each composition $X_i \to Y$ has $\mathcal{P}$,
\item[(g)] there exists a smooth covering $\{Y_i \to Y\}_{i \in I}$ and
for each $i \in I$ a smooth covering
$\{X_{ij} \to Y_i \times_Y X\}_{j \in J_i}$ such that each morphism
$X_{ij} \to Y_i$ has $\mathcal{P}$.
\end{enumerate}
\end{lemma}

\begin{proof}
The equivalence of (a) and (b) is part of
Definition \ref{definition-local-source-target}.
The equivalence of (a) and (e) is
Lemma \ref{lemma-local-source-target-implies} part (2).
The equivalence of (a) and (f) is
Lemma \ref{lemma-local-source-target-implies} part (1).
As (a) is now equivalent to (e) and (f) it follows that
(a) equivalent to (g).

\medskip\noindent
It is clear that (c) implies (b). If (b) holds, then for any
$x \in |X|$ we can choose a smooth morphism of pairs
$a_x : (U_x, u_x) \to (X, x)$, a smooth morphism $b_x : V_x \to Y$, and
a morphism $h_x : U_x \to V_x$ such that $f \circ a_x = b_x \circ h_x$ and
$h_x$ has $\mathcal{P}$. Then $h = \coprod h_x : \coprod U_x \to \coprod V_x$
with $a = \coprod a_x$ and $b = \coprod b_x$ is a diagram as in (c).
(Note that $h$ has property $\mathcal{P}$ as $\{V_x \to \coprod V_x\}$
is a smooth covering and $\mathcal{P}$ is smooth local on the target.)
Thus (b) is equivalent to (c).

\medskip\noindent
Now we know that (a), (b), (c), (e), (f), and (g) are equivalent.
Suppose (a) holds. Let $U, V, a, b, h$ be as in (d). Then
$X \times_Y V \to V$ has $\mathcal{P}$ as $\mathcal{P}$ is stable under
smooth base change, whence $U \to V$ has $\mathcal{P}$ as $\mathcal{P}$
is stable under precomposing with smooth morphisms. Conversely, if (d)
holds, then setting $U = X$ and $V = Y$ we see that $f$ has $\mathcal{P}$.
\end{proof}

\begin{lemma}
\label{lemma-smooth-local-source-target}
Let $S$ be a scheme.
Let $\mathcal{P}$ be a property of morphisms of algebraic spaces over $S$.
Assume
\begin{enumerate}
\item $\mathcal{P}$ is smooth local on the source,
\item $\mathcal{P}$ is smooth local on the target, and
\item $\mathcal{P}$ is stable under postcomposing with smooth morphisms:
if $f : X \to Y$ has $\mathcal{P}$ and $Y \to Z$ is a smooth morphism
then $X \to Z$ has $\mathcal{P}$.
\end{enumerate}
Then $\mathcal{P}$ is smooth local on the source-and-target.
\end{lemma}

\begin{proof}
Let $\mathcal{P}$ be a property of morphisms of algebraic spaces which
satisfies conditions (1), (2) and (3) of the lemma. By
Lemma \ref{lemma-precompose-property-local-source}
we see that $\mathcal{P}$ is stable under precomposing with
smooth morphisms. By
Lemma \ref{lemma-pullback-property-local-target}
we see that $\mathcal{P}$ is stable under smooth base change.
Hence it suffices to prove part (3) of
Definition \ref{definition-local-source-target}
holds.

\medskip\noindent
More precisely, suppose that $f : X \to Y$ is a morphism
of algebraic spaces over $S$ which satisfies
Definition \ref{definition-local-source-target} part (3)(b).
In other words, for every $x \in X$ there exists a smooth
morphism $a_x : U_x \to X$, a point $u_x \in |U_x|$ mapping to $x$,
a smooth morphism $b_x : V_x \to Y$, and a morphism $h_x : U_x \to V_x$
such that $f \circ a_x = b_x \circ h_x$ and $h_x$ has $\mathcal{P}$.
The proof of the lemma is complete once we show that $f$ has $\mathcal{P}$.
Set $U = \coprod U_x$, $a = \coprod a_x$, $V = \coprod V_x$,
$b = \coprod b_x$, and $h = \coprod h_x$. We obtain a
commutative diagram
$$
\xymatrix{
U \ar[d]_a \ar[r]_h & V \ar[d]^b \\
X \ar[r]^f & Y
}
$$
with $a$, $b$ smooth, $a$ surjective. Note that $h$ has $\mathcal{P}$
as each $h_x$ does and $\mathcal{P}$ is smooth local on the target.
Because $a$ is surjective and $\mathcal{P}$ is smooth local on the source,
it suffices to prove that $b \circ h$ has $\mathcal{P}$.
This follows as we assumed that $\mathcal{P}$ is stable under
postcomposing with a smooth morphism and as $b$ is smooth.
\end{proof}

\begin{remark}
\label{remark-list-local-source-target}
Using
Lemma \ref{lemma-smooth-local-source-target}
and the work done in the earlier sections of this chapter it is easy
to make a list of types of morphisms which are smooth local on the
source-and-target. In each case we list the lemma which implies
the property is smooth local on the source and the lemma which implies
the property is smooth local on the target. In each case the third assumption
of
Lemma \ref{lemma-smooth-local-source-target}
is trivial to check, and we omit it. Here is the list:
\begin{enumerate}
\item flat, see
Lemmas \ref{lemma-flat-fpqc-local-source} and
\ref{lemma-descending-property-flat},
\item locally of finite presentation, see
Lemmas \ref{lemma-locally-finite-presentation-fppf-local-source} and
\ref{lemma-descending-property-locally-finite-presentation},
\item locally finite type, see
Lemmas \ref{lemma-locally-finite-type-fppf-local-source} and
\ref{lemma-descending-property-locally-finite-type},
\item universally open, see
Lemmas \ref{lemma-universally-open-fppf-local-source} and
\ref{lemma-descending-property-universally-open},
\item syntomic, see
Lemmas \ref{lemma-syntomic-syntomic-local-source} and
\ref{lemma-descending-property-syntomic},
\item smooth, see
Lemmas \ref{lemma-smooth-smooth-local-source} and
\ref{lemma-descending-property-smooth},
\item add more here as needed.
\end{enumerate}
\end{remark}







\section{Properties of morphisms \'etale-smooth local on source-and-target}
\label{section-properties-etale-smooth-local-source-target}

\noindent
This section is the analogue of
Section \ref{section-properties-local-source-target}
for properties of morphisms which are \'etale local
on the source and smooth local on the target.
We give this property a ridiculously long name
in order to avoid using it too much.

\begin{definition}
\label{definition-etale-smooth-local-source-target}
Let $S$ be a scheme.
Let $\mathcal{P}$ be a property of morphisms of algebraic spaces over $S$.
We say $\mathcal{P}$ is {\it \'etale-smooth local on source-and-target} if
\begin{enumerate}
\item (stable under precomposing with \'etale maps)
if $f : X \to Y$ is \'etale and $g : Y \to Z$ has $\mathcal{P}$,
then $g \circ f$ has $\mathcal{P}$,
\item (stable under smooth base change)
if $f : X \to Y$ has $\mathcal{P}$ and $Y' \to Y$ is smooth, then
the base change $f' : Y' \times_Y X \to Y'$ has $\mathcal{P}$, and
\item (locality) given a morphism $f : X \to Y$ the following are
equivalent
\begin{enumerate}
\item $f$ has $\mathcal{P}$,
\item for every $x \in |X|$ there exists a commutative diagram
$$
\xymatrix{
U \ar[d]_a \ar[r]_h & V \ar[d]^b \\
X \ar[r]^f & Y
}
$$
with $b$ smooth and $U \to X \times_Y V$ \'etale
and $u \in |U|$ with $a(u) = x$ such that
$h$ has $\mathcal{P}$.
\end{enumerate}
\end{enumerate}
\end{definition}

\noindent
The above serves as our definition. In the lemmas below we will show that
this is equivalent to $\mathcal{P}$ being \'etale local on the target,
smooth local on the source, and stable under post-composing by
\'etale morphisms.

\begin{lemma}
\label{lemma-etale-smooth-local-source-target-implies}
Let $S$ be a scheme.
Let $\mathcal{P}$ be a property of morphisms of algebraic spaces over $S$
which is \'etale-smooth local on source-and-target. Then
\begin{enumerate}
\item $\mathcal{P}$ is \'etale local on the source,
\item $\mathcal{P}$ is smooth local on the target,
\item $\mathcal{P}$ is stable under postcomposing with \'etale morphisms:
if $f : X \to Y$ has $\mathcal{P}$ and $g : Y \to Z$ is \'etale, then
$g \circ f$ has $\mathcal{P}$, and
\item $\mathcal{P}$ has a permanence property: given $f : X \to Y$ and
$g : Y \to Z$ \'etale such that $g \circ f$ has $\mathcal{P}$, then
$f$ has $\mathcal{P}$.
\end{enumerate}
\end{lemma}

\begin{proof}
We write everything out completely.

\medskip\noindent
Proof of (1). Let $f : X \to Y$ be a morphism of algebraic spaces over $S$.
Let $\{X_i \to X\}_{i \in I}$ be an \'etale covering of $X$. If each composition
$h_i : X_i \to Y$ has $\mathcal{P}$, then for each $|x| \in X$ we can find
an $i \in I$ and a point $x_i \in |X_i|$ mapping to $x$. Then
$(X_i, x_i) \to (X, x)$ is an \'etale morphism of pairs, and
$\text{id}_Y : Y \to Y$ is a smooth morphism, and $h_i$ is as in part (3) of
Definition \ref{definition-etale-smooth-local-source-target}.
Thus we see that $f$ has $\mathcal{P}$.
Conversely, if $f$ has $\mathcal{P}$ then each $X_i \to Y$ has
$\mathcal{P}$ by
Definition \ref{definition-etale-smooth-local-source-target} part (1).

\medskip\noindent
Proof of (2). Let $f : X \to Y$ be a morphism of algebraic spaces over $S$.
Let $\{Y_i \to Y\}_{i \in I}$ be a smooth covering of $Y$.
Write $X_i = Y_i \times_Y X$ and $h_i : X_i \to Y_i$ for the base change
of $f$. If each  $h_i : X_i \to Y_i$ has $\mathcal{P}$, then for each
$x \in |X|$ we pick an $i \in I$ and a point $x_i \in |X_i|$ mapping to $x$.
Then $X_i \to X \times_Y Y_i$ is an \'etale morphism
(because it is an isomorphism), $Y_i \to Y$ is
smooth, and $h_i$ is as in part (3) of
Definition \ref{definition-local-source-target}.
Thus we see that $f$ has $\mathcal{P}$.
Conversely, if $f$ has $\mathcal{P}$, then each $X_i \to Y_i$ has
$\mathcal{P}$ by
Definition \ref{definition-local-source-target} part (2).

\medskip\noindent
Proof of (3). Assume $f : X \to Y$ has $\mathcal{P}$ and $g : Y \to Z$ is
\'etale. The morphism $X \to Y \times_Z X$ is \'etale as a morphism
between algebraic spaces \'etale over $X$ (
Properties of Spaces, Lemma \ref{spaces-properties-lemma-etale-permanence}).
Also $Y \to Z$ is \'etale hence a smooth morphism.
Thus the diagram
$$
\xymatrix{
X \ar[d] \ar[r]_f & Y \ar[d] \\
X \ar[r]^{g \circ f} & Z
}
$$
works for every $x \in |X|$ in part (3) of
Definition \ref{definition-local-source-target}
and we conclude that $g \circ f$ has $\mathcal{P}$.

\medskip\noindent
Proof of (4). Let $f : X \to Y$ be a morphism and $g : Y \to Z$ \'etale
such that $g \circ f$ has $\mathcal{P}$. Then by
Definition \ref{definition-etale-smooth-local-source-target} part (2)
we see that $\text{pr}_Y : Y \times_Z X \to Y$ has $\mathcal{P}$. But
the morphism $(f, 1) : X \to Y \times_Z X$ is \'etale as a section to the
\'etale projection $\text{pr}_X : Y \times_Z X \to X$, see
Morphisms of Spaces, Lemma \ref{spaces-morphisms-lemma-etale-permanence}.
Hence $f = \text{pr}_Y \circ (f, 1)$ has $\mathcal{P}$ by
Definition \ref{definition-etale-smooth-local-source-target} part (1).
\end{proof}

\begin{lemma}
\label{lemma-etale-smooth-local-source-target-characterize}
Let $S$ be a scheme.
Let $\mathcal{P}$ be a property of morphisms of algebraic spaces over $S$
which is etale-smooth local on source-and-target.
Let $f : X \to Y$ be a morphism
of algebraic spaces over $S$. The following are equivalent:
\begin{enumerate}
\item[(a)] $f$ has property $\mathcal{P}$,
\item[(b)] for every $x \in |X|$ there exists a smooth morphism $b : V \to Y$,
an \'etale morphism $a : U \to V \times_Y X$, and a point $u \in |U|$
mapping to $x$ such that $U \to V$ has $\mathcal{P}$,
\item[(c)] for some commutative diagram
$$
\xymatrix{
U \ar[d]_a \ar[r]_h & V \ar[d]^b \\
X \ar[r]^f & Y
}
$$
with $b$ smooth, $U \to V \times_Y X$ \'etale, and $a$ surjective
the morphism $h$ has $\mathcal{P}$,
\item[(d)] for any commutative diagram
$$
\xymatrix{
U \ar[d]_a \ar[r]_h & V \ar[d]^b \\
X \ar[r]^f & Y
}
$$
with $b$ smooth and $U \to X \times_Y V$ \'etale, the morphism $h$
has $\mathcal{P}$,
\item[(e)] there exists a smooth covering $\{Y_i \to Y\}_{i \in I}$ such
that each base change $Y_i \times_Y X \to Y_i$ has $\mathcal{P}$,
\item[(f)] there exists an \'etale covering $\{X_i \to X\}_{i \in I}$ such
that each composition $X_i \to Y$ has $\mathcal{P}$,
\item[(g)] there exists a smooth covering $\{Y_i \to Y\}_{i \in I}$ and
for each $i \in I$ an \'etale covering
$\{X_{ij} \to Y_i \times_Y X\}_{j \in J_i}$ such that each morphism
$X_{ij} \to Y_i$ has $\mathcal{P}$.
\end{enumerate}
\end{lemma}

\begin{proof}
The equivalence of (a) and (b) is part of
Definition \ref{definition-etale-smooth-local-source-target}.
The equivalence of (a) and (e) is
Lemma \ref{lemma-etale-smooth-local-source-target-implies} part (2).
The equivalence of (a) and (f) is
Lemma \ref{lemma-etale-smooth-local-source-target-implies} part (1).
As (a) is now equivalent to (e) and (f) it follows that
(a) equivalent to (g).

\medskip\noindent
It is clear that (c) implies (b). If (b) holds, then for any
$x \in |X|$ we can choose a smooth morphism a smooth morphism
$b_x : V_x \to Y$, an \'etale morphism $U_x \to V_x \times_Y X$,
and $u_x \in |U_x|$ mapping to $x$
such that $U_x \to V_x$ has $\mathcal{P}$.
Then $h = \coprod h_x : \coprod U_x \to \coprod V_x$
with $a = \coprod a_x$ and $b = \coprod b_x$ is a diagram as in (c).
(Note that $h$ has property $\mathcal{P}$ as $\{V_x \to \coprod V_x\}$
is a smooth covering and $\mathcal{P}$ is smooth local on the target.)
Thus (b) is equivalent to (c).

\medskip\noindent
Now we know that (a), (b), (c), (e), (f), and (g) are equivalent.
Suppose (a) holds. Let $U, V, a, b, h$ be as in (d). Then
$X \times_Y V \to V$ has $\mathcal{P}$ as $\mathcal{P}$ is stable under
smooth base change, whence $U \to V$ has $\mathcal{P}$ as $\mathcal{P}$
is stable under precomposing with \'etale morphisms. Conversely, if (d)
holds, then setting $U = X$ and $V = Y$ we see that $f$ has $\mathcal{P}$.
\end{proof}

\begin{lemma}
\label{lemma-etale-smooth-local-source-target}
Let $S$ be a scheme.
Let $\mathcal{P}$ be a property of morphisms of algebraic spaces over $S$.
Assume
\begin{enumerate}
\item $\mathcal{P}$ is \'etale local on the source,
\item $\mathcal{P}$ is smooth local on the target, and
\item $\mathcal{P}$ is stable under postcomposing with open immersions:
if $f : X \to Y$ has $\mathcal{P}$ and $Y \subset Z$ is an open embedding
then $X \to Z$ has $\mathcal{P}$.
\end{enumerate}
Then $\mathcal{P}$ is \'etale-smooth local on the source-and-target.
\end{lemma}

\begin{proof}
Let $\mathcal{P}$ be a property of morphisms of algebraic spaces which
satisfies conditions (1), (2) and (3) of the lemma. By
Lemma \ref{lemma-precompose-property-local-source}
we see that $\mathcal{P}$ is stable under precomposing with
\'etale morphisms. By
Lemma \ref{lemma-pullback-property-local-target}
we see that $\mathcal{P}$ is stable under smooth base change.
Hence it suffices to prove part (3) of
Definition \ref{definition-local-source-target}
holds.

\medskip\noindent
More precisely, suppose that $f : X \to Y$ is a morphism
of algebraic spaces over $S$ which satisfies
Definition \ref{definition-local-source-target} part (3)(b).
In other words, for every $x \in X$ there exists
a smooth morphism $b_x : V_x \to Y$,
an \'etale morphism $U_x \to V_x \times_Y X$, and
a point $u_x \in |U_x|$ mapping to $x$
such that $h_x : U_x \to V_x$ has $\mathcal{P}$.
The proof of the lemma is complete once we show that $f$ has $\mathcal{P}$.

\medskip\noindent
Let $a_x : U_x \to X$ be the composition $U_x \to V_x \times_Y X \to X$.
Set $U = \coprod U_x$, $a = \coprod a_x$, $V = \coprod V_x$,
$b = \coprod b_x$, and $h = \coprod h_x$. We obtain a
commutative diagram
$$
\xymatrix{
U \ar[d]_a \ar[r]_h & V \ar[d]^b \\
X \ar[r]^f & Y
}
$$
with $b$ smooth, $U \to V \times_Y X$ \'etale, $a$ surjective.
Note that $h$ has $\mathcal{P}$ as each $h_x$ does and $\mathcal{P}$
is smooth local on the target. In the next paragraph we prove
that we may assume $U, V, X, Y$ are schemes; we encourage the reader
to skip it.

\medskip\noindent
Let $X, Y, U, V, a, b, f, h$ be as in the previous paragraph. We have
to show $f$ has $\mathcal{P}$. Let $X' \to X$ be a surjective \'etale
morphism with $X_i$ a scheme. Set $U' = X' \times_X U$. Then
$U' \to X'$ is surjective and $U' \to X' \times_Y V$ is \'etale.
Since $\mathcal{P}$ is \'etale local on the source, we see that
$U' \to V$ has $\mathcal{P}$ and that it suffices to show that
$X' \to Y$ has $\mathcal{P}$. In other words, we may assume
that $X$ is a scheme. Next, choose a surjective \'etale morphism
$Y' \to Y$ with $Y'$ a scheme. Set $V' = V \times_Y Y'$,
$X' = X \times_Y Y'$, and $U' = U \times_Y Y'$. Then
$U' \to X'$ is surjective and $U' \to X' \times_{Y'} V'$ is \'etale.
Since $\mathcal{P}$ is smooth local on the target, we see that $U' \to V'$ has
$\mathcal{P}$ and that it suffices to prove $X' \to Y'$ has $\mathcal{P}$.
Thus we may assume both $X$ and $Y$ are schemes.
Choose a surjective \'etale morphism $V' \to V$
with $V'$ a scheme. Set $U' = U \times_V V'$.
Then $U' \to X$ is surjective and $U' \to X \times_Y V'$ is \'etale.
Since $\mathcal{P}$ is smooth local on the source, we see that
$U' \to V'$ has $\mathcal{P}$. Thus we may replace $U, V$ by
$U', V'$ and assume $X, Y, V$ are schemes.
Finally, we replace $U$ by a scheme surjective \'etale over $U$
and we see that we may assume $U, V, X, Y$ are all schemes.

\medskip\noindent
If $U, V, X, Y$ are schemes, then $f$ has $\mathcal{P}$
by Descent, Lemma \ref{descent-lemma-etale-tau-local-source-target}.
\end{proof}

\begin{remark}
\label{remark-list-etale-smooth-local-source-target}
Using Lemma \ref{lemma-etale-smooth-local-source-target}
and the work done in the earlier sections of this chapter it is easy
to make a list of types of morphisms which are smooth local on the
source-and-target. In each case we list the lemma which implies
the property is etale local on the source and the lemma which implies
the property is smooth local on the target. In each case the third assumption
of Lemma \ref{lemma-etale-smooth-local-source-target}
is trivial to check, and we omit it. Here is the list:
\begin{enumerate}
\item \'etale, see
Lemmas \ref{lemma-etale-etale-local-source} and
\ref{lemma-descending-property-etale},
\item locally quasi-finite, see
Lemmas \ref{lemma-locally-quasi-finite-etale-local-source} and
\ref{lemma-descending-property-quasi-finite},
\item unramified, see
Lemmas \ref{lemma-unramified-etale-local-source} and
\ref{lemma-descending-property-unramified}, and
\item add more here as needed.
\end{enumerate}
Of course any property listed in
Remark \ref{remark-list-local-source-target}
is a fortiori an example that could be listed here.
\end{remark}







\section{Descent data for spaces over spaces}
\label{section-descent-datum}

\noindent
This section is the analogue of Descent, Section
\ref{descent-section-descent-datum} for algebraic spaces.
Most of the arguments in this section are formal relying only
on the definition of a descent datum.

\begin{definition}
\label{definition-descent-datum}
Let $S$ be a scheme. Let $f : Y \to X$ be a morphism of algebraic spaces
over $S$.
\begin{enumerate}
\item Let $V \to Y$ be a morphism of algebraic spaces.
A {\it descent datum for $V/Y/X$} is an isomorphism
$\varphi : V \times_X Y \to Y \times_X V$ of algebraic spaces over
$Y \times_X Y$ satisfying the {\it cocycle condition} that the diagram
$$
\xymatrix{
V \times_X Y \times_X Y \ar[rd]^{\varphi_{01}} \ar[rr]_{\varphi_{02}} &
&
Y \times_X Y \times_X V\\
&
Y \times_X V \times_X Y \ar[ru]^{\varphi_{12}}
}
$$
commutes (with obvious notation).
\item We also say that the pair $(V/Y, \varphi)$ is
a {\it descent datum relative to $Y \to X$}.
\item A {\it morphism $f : (V/Y, \varphi) \to (V'/Y, \varphi')$ of
descent data relative to $Y \to X$} is a morphism
$f : V \to V'$ of algebraic spaces over $Y$ such that
the diagram
$$
\xymatrix{
V \times_X Y \ar[r]_{\varphi} \ar[d]_{f \times \text{id}_Y} &
Y \times_X V \ar[d]^{\text{id}_Y \times f} \\
V' \times_X Y \ar[r]^{\varphi'} & Y \times_X V'
}
$$
commutes.
\end{enumerate}
\end{definition}

\begin{remark}
\label{remark-easier}
Let $S$ be a scheme.
Let $Y \to X$ be a morphism of algebraic spaces over $S$.
Let $(V/Y, \varphi)$ be a descent datum relative to $Y \to X$.
We may think of the isomorphism $\varphi$ as an isomorphism
$$
(Y \times_X Y) \times_{\text{pr}_0, Y} V
\longrightarrow
(Y \times_X Y) \times_{\text{pr}_1, Y} V
$$
of algebraic spaces over $Y \times_X Y$. So loosely speaking one may
think of $\varphi$ as a map
$\varphi : \text{pr}_0^*V \to \text{pr}_1^*V$\footnote{Unfortunately,
we have chosen the ``wrong'' direction for our arrow here. In
Definitions \ref{definition-descent-datum} and
\ref{definition-descent-datum-for-family-of-morphisms}
we should have the opposite direction to what was done in
Definition \ref{definition-descent-datum-quasi-coherent}
by the general principle that ``functions'' and ``spaces'' are dual.}.
The cocycle condition then says that
$\text{pr}_{02}^*\varphi =
\text{pr}_{12}^*\varphi \circ \text{pr}_{01}^*\varphi$.
In this way it is very similar to the case of a descent datum on
quasi-coherent sheaves.
\end{remark}

\noindent
Here is the definition in case you have a family of morphisms
with fixed target.

\begin{definition}
\label{definition-descent-datum-for-family-of-morphisms}
Let $S$ be a scheme.
Let $\{X_i \to X\}_{i \in I}$ be a family of morphisms
of algebraic spaces over $S$ with fixed target $X$.
\begin{enumerate}
\item A {\it descent datum $(V_i, \varphi_{ij})$ relative to the
family $\{X_i \to X\}$} is given by an algebraic space $V_i$ over $X_i$
for each $i \in I$, an isomorphism
$\varphi_{ij} : V_i \times_X X_j \to X_i \times_X V_j$
of algebraic spaces over $X_i \times_X X_j$ for each pair $(i, j) \in I^2$
such that for every triple of indices $(i, j, k) \in I^3$
the diagram
$$
\xymatrix{
V_i \times_X X_j \times_X X_k
\ar[rd]^{\text{pr}_{01}^*\varphi_{ij}}
\ar[rr]_{\text{pr}_{02}^*\varphi_{ik}} &
&
X_i \times_X X_j \times_X V_k\\
&
X_i \times_X V_j \times_X X_k
\ar[ru]^{\text{pr}_{12}^*\varphi_{jk}}
}
$$
of algebraic spaces over $X_i \times_X X_j \times_X X_k$ commutes
(with obvious notation).
\item A {\it morphism
$\psi : (V_i, \varphi_{ij}) \to (V'_i, \varphi'_{ij})$
of descent data} is given by a family $\psi = (\psi_i)_{i \in I}$
of morphisms $\psi_i : V_i \to V'_i$ of algebraic spaces over $X_i$
such that all the diagrams
$$
\xymatrix{
V_i \times_X X_j \ar[r]_{\varphi_{ij}} \ar[d]_{\psi_i \times \text{id}} &
X_i \times_X V_j \ar[d]^{\text{id} \times \psi_j} \\
V'_i \times_X X_j \ar[r]^{\varphi'_{ij}} & X_i \times_X V'_j
}
$$
commute.
\end{enumerate}
\end{definition}

\begin{remark}
\label{remark-easier-family}
Let $S$ be a scheme.
Let $\{X_i \to X\}_{i \in I}$ be a family of morphisms
of algebraic spaces over $S$ with fixed target $X$.
Let $(V_i, \varphi_{ij})$ be a descent datum relative to
$\{X_i \to X\}$. We may think of the isomorphisms $\varphi_{ij}$
as isomorphisms
$$
(X_i \times_X X_j) \times_{\text{pr}_0, X_i} V_i
\longrightarrow
(X_i \times_X X_j) \times_{\text{pr}_1, X_j} V_j
$$
of algebraic spaces over $X_i \times_X X_j$. So loosely speaking one may
think of $\varphi_{ij}$ as an isomorphism
$\text{pr}_0^*V_i \to \text{pr}_1^*V_j$ over $X_i \times_X X_j$.
The cocycle condition then says that
$\text{pr}_{02}^*\varphi_{ik} =
\text{pr}_{12}^*\varphi_{jk} \circ \text{pr}_{01}^*\varphi_{ij}$.
In this way it is very similar to the case of a descent datum on
quasi-coherent sheaves.
\end{remark}

\noindent
The reason we will usually work with the version of a family consisting
of a single morphism is the following lemma.

\begin{lemma}
\label{lemma-family-is-one}
Let $S$ be a scheme.
Let $\{X_i \to X\}_{i \in I}$ be a family of morphisms
of algebraic spaces over $S$ with fixed target $X$.
Set $Y = \coprod_{i \in I} X_i$.
There is a canonical equivalence of categories
$$
\begin{matrix}
\text{category of descent data } \\
\text{relative to the family } \{X_i \to X\}_{i \in I}
\end{matrix}
\longrightarrow
\begin{matrix}
\text{ category of descent data} \\
\text{ relative to } Y/X
\end{matrix}
$$
which maps $(V_i, \varphi_{ij})$ to $(V, \varphi)$ with
$V = \coprod_{i\in I} V_i$ and $\varphi = \coprod \varphi_{ij}$.
\end{lemma}

\begin{proof}
Observe that $Y \times_X Y = \coprod_{ij} X_i \times_X X_j$
and similarly for higher fibre products.
Giving a morphism $V \to Y$ is exactly the same as
giving a family $V_i \to X_i$. And giving a descent datum
$\varphi$ is exactly the same as giving a family $\varphi_{ij}$.
\end{proof}

\begin{lemma}
\label{lemma-pullback}
Pullback of descent data. Let $S$ be a scheme.
\begin{enumerate}
\item Let
$$
\xymatrix{
Y' \ar[r]_f \ar[d]_{a'} & Y \ar[d]^a \\
X' \ar[r]^h & X
}
$$
be a commutative diagram of algebraic spaces over $S$.
The construction
$$
(V \to Y, \varphi) \longmapsto f^*(V \to Y, \varphi) = (V' \to Y', \varphi')
$$
where $V' = Y' \times_Y V$ and where
$\varphi'$ is defined as the composition
$$
\xymatrix{
V' \times_{X'} Y' \ar@{=}[r] &
(Y' \times_Y V) \times_{X'} Y' \ar@{=}[r] &
(Y' \times_{X'} Y') \times_{Y \times_X Y} (V \times_X Y)
\ar[d]^{\text{id} \times \varphi} \\
Y' \times_{X'} V' \ar@{=}[r] &
Y' \times_{X'} (Y' \times_Y V) &
(Y' \times_X Y') \times_{Y \times_X Y} (Y \times_X V) \ar@{=}[l]
}
$$
defines a functor from the category of descent data
relative to $Y \to X$ to the category of descent data
relative to $Y' \to X'$.
\item Given two morphisms $f_i : Y' \to Y$, $i = 0, 1$ making the
diagram commute the functors $f_0^*$ and $f_1^*$ are
canonically isomorphic.
\end{enumerate}
\end{lemma}

\begin{proof}
We omit the proof of (1), but we remark that the morphism
$\varphi'$ is the morphism $(f \times f)^*\varphi$ in the
notation introduced in Remark \ref{remark-easier}.
For (2) we indicate which morphism
$f_0^*V \to f_1^*V$ gives the functorial isomorphism. Namely,
since $f_0$ and $f_1$ both fit into the commutative diagram
we see there is a unique morphism $r : Y' \to Y \times_X Y$
with $f_i = \text{pr}_i \circ r$. Then we take
\begin{eqnarray*}
f_0^*V & = &
Y' \times_{f_0, Y} V \\
& = &
Y' \times_{\text{pr}_0 \circ r, Y} V \\
& = &
Y' \times_{r, Y \times_X Y} (Y \times_X Y) \times_{\text{pr}_0, Y} V \\
& \xrightarrow{\varphi} &
Y' \times_{r, Y \times_X Y} (Y \times_X Y) \times_{\text{pr}_1, Y} V \\
& = &
Y' \times_{\text{pr}_1 \circ r, Y} V \\
& = &
Y' \times_{f_1, Y} V \\
& = & f_1^*V
\end{eqnarray*}
We omit the verification that this works.
\end{proof}

\begin{definition}
\label{definition-pullback-functor}
With $S, X, X', Y, Y', f, a, a', h$ as in Lemma \ref{lemma-pullback}
the functor
$$
(V, \varphi) \longmapsto f^*(V, \varphi)
$$
constructed in that lemma is called the {\it pullback functor} on descent data.
\end{definition}

\begin{lemma}
\label{lemma-pullback-family}
Let $S$ be a scheme. Let $\mathcal{U}' = \{X'_i \to X'\}_{i \in I'}$ and
$\mathcal{U} = \{X_j \to X\}_{i \in I}$ be families of morphisms with
fixed target. Let $\alpha : I' \to I$, $g : X' \to X$ and
$g_i : X'_i \to X_{\alpha(i)}$ be a morphism of families
of maps with fixed target, see
Sites, Definition \ref{sites-definition-morphism-coverings}.
\begin{enumerate}
\item Let $(V_i, \varphi_{ij})$ be a descent datum relative to the
family $\mathcal{U}$. The system
$$
\left(
g_i^*V_{\alpha(i)}, (g_i \times g_j)^*\varphi_{\alpha(i) \alpha(j)}
\right)
$$
(with notation as in Remark \ref{remark-easier-family})
is a descent datum relative to $\mathcal{U}'$.
\item This construction defines a functor between the category of
descent data relative to $\mathcal{U}$ and the category of
descent data relative to $\mathcal{U}'$.
\item Given a second $\beta : I' \to I$, $h : X' \to X$ and
$h'_i : X'_i \to X_{\beta(i)}$ morphism of families
of maps with fixed target, then if $g = h$ the two resulting functors
between descent data are canonically isomorphic.
\item These functors agree, via Lemma \ref{lemma-family-is-one},
with the pullback functors constructed in Lemma \ref{lemma-pullback}.
\end{enumerate}
\end{lemma}

\begin{proof}
This follows from Lemma \ref{lemma-pullback} via the
correspondence of Lemma \ref{lemma-family-is-one}.
\end{proof}

\begin{definition}
\label{definition-pullback-functor-family}
With $\mathcal{U}' = \{X'_i \to X'\}_{i \in I'}$,
$\mathcal{U} = \{X_i \to X\}_{i \in I}$, $\alpha : I' \to I$,
$g : X' \to X$, and $g_i : X'_i \to X_{\alpha(i)}$ as in
Lemma \ref{lemma-pullback-family} the functor
$$
(V_i, \varphi_{ij}) \longmapsto
(g_i^*V_{\alpha(i)}, (g_i \times g_j)^*\varphi_{\alpha(i) \alpha(j)})
$$
constructed in that lemma
is called the {\it pullback functor} on descent data.
\end{definition}

\noindent
If $\mathcal{U}$ and $\mathcal{U}'$ have the same target $X$,
and if $\mathcal{U}'$ refines $\mathcal{U}$ (see
Sites, Definition \ref{sites-definition-morphism-coverings})
but no explicit pair $(\alpha, g_i)$ is given, then we can still
talk about the pullback functor since we have seen in
Lemma \ref{lemma-pullback-family} that the choice of the pair does not matter
(up to a canonical isomorphism).

\begin{definition}
\label{definition-effective}
Let $S$ be a scheme. Let $f : Y \to X$ be a morphism of algebraic spaces over
$S$.
\begin{enumerate}
\item Given an algebraic space $U$ over $X$ we have the
{\it trivial descent datum} of $U$ relative to $\text{id} : X \to X$, namely
the identity morphism on $U$.
\item By Lemma \ref{lemma-pullback} we get a
{\it canonical descent datum} on $Y \times_X U$
relative to $Y \to X$ by pulling back the trivial
descent datum via $f$. We often
denote $(Y \times_X U, can)$ this descent datum.
\item A descent datum $(V, \varphi)$ relative to $Y/X$
is called {\it effective} if $(V, \varphi)$
is isomorphic to the canonical descent datum
$(Y \times_X U, can)$ for some algebraic space $U$ over $X$.
\end{enumerate}
\end{definition}

\noindent
Thus being effective means there exists an algebraic space $U$
over $X$ and an isomorphism $\psi : V \to Y \times_X U$
over $Y$ such that $\varphi$ is equal to the composition
$$
V \times_X Y \xrightarrow{\psi \times \text{id}_Y}
Y \times_X U \times_S Y =
Y \times_X Y \times_X U
\xrightarrow{\text{id}_Y \times \psi^{-1}}
Y \times_X V
$$
There is a slight problem here which is that this definition
(in spirit) conflicts with the definition given in
Descent, Definition \ref{descent-definition-effective}
in case $Y$ and $X$ are schemes. However, it will always be clear from
context which version we mean.

\begin{definition}
\label{definition-effective-family}
Let $S$ be a scheme.
Let $\{X_i \to X\}$ be a family of morphisms of algebraic spaces over $S$
with fixed target $X$.
\begin{enumerate}
\item  Given an algebraic space $U$ over $X$
we have a {\it canonical descent datum} on the family of
algebraic spaces $X_i \times_X U$ by pulling back the trivial
descent datum for $U$ relative to $\{\text{id} : S \to S\}$.
We denote this descent datum $(X_i \times_X U, can)$.
\item A descent datum $(V_i, \varphi_{ij})$
relative to $\{X_i \to S\}$ is called {\it effective}
if there exists an algebraic space $U$ over $X$ such that
$(V_i, \varphi_{ij})$ is isomorphic to $(X_i \times_X U, can)$.
\end{enumerate}
\end{definition}






\section{Descent data in terms of sheaves}
\label{section-descent-data-sheaves}

\noindent
This section is the analogue of Descent, Section
\ref{descent-section-descent-data-sheaves}.
It is slightly different as algebraic spaces are already sheaves.

\begin{lemma}
\label{lemma-descent-data-sheaves}
Let $S$ be a scheme. Let $\{X_i \to X\}_{i \in I}$ be an fppf
covering of algebraic spaces over $S$ (Topologies on Spaces,
Definition \ref{spaces-topologies-definition-fppf-covering}).
There is an equivalence of categories
$$
\left\{
\begin{matrix}
\text{descent data }(V_i, \varphi_{ij})\\
\text{relative to }\{X_i \to X\}
\end{matrix}
\right\}
\leftrightarrow
\left\{
\begin{matrix}
\text{sheaves }F\text{ on }(\Sch/S)_{fppf}\text{ endowed}\\
\text{with a map }F \to X\text{ such that each}\\
X_i \times_X F\text{ is an algebraic space}
\end{matrix}
\right\}.
$$
Moreover,
\begin{enumerate}
\item the algebraic space $X_i \times_X F$ on the right hand side
corresponds to $V_i$ on the left hand side, and
\item the sheaf $F$ is an algebraic space\footnote{We will see
later that this is always the case if $I$ is not too large, see
Bootstrap, Lemma \ref{bootstrap-lemma-descend-algebraic-space}.}
if and only if the
corresponding descent datum $(X_i, \varphi_{ij})$ is effective.
\end{enumerate}
\end{lemma}

\begin{proof}
Let us construct the functor from right to left.
Let $F \to X$ be a map of sheaves on $(\Sch/S)_{fppf}$ such that each
$V_i = X_i \times_X F$ is an algebraic space. We have the
projection $V_i \to X_i$.
Then both $V_i \times_X X_j$ and $X_i \times_X V_j$
represent the sheaf $X_i \times_X F \times_X X_j$
and hence we obtain an isomorphism
$$
\varphi_{ii'} : V_i \times_X X_j \to X_i \times_X V_j
$$
It is straightforward to see that the maps $\varphi_{ij}$
are morphisms over $X_i \times_X X_j$ and satisfy the
cocycle condition. The functor from right to left is given
by this construction $F \mapsto (V_i, \varphi_{ij})$.

\medskip\noindent
Let us construct a functor from left to right.
The isomorphisms $\varphi_{ij}$ give isomorphisms
$$
\varphi_{ij} : V_i \times_X X_j \longrightarrow X_i \times_X V_j
$$
over $X_i \times X_j$. Set $F$ equal to the coequalizer in the
following diagram
$$
\xymatrix{
\coprod_{i, i'} V_i \times_X X_j
\ar@<1ex>[rr]^-{\text{pr}_0}
\ar@<-1ex>[rr]_-{\text{pr}_1 \circ \varphi_{ij}}
& &
\coprod_i V_i \ar[r]
&
F
}
$$
The cocycle condition guarantees that $F$ comes with a map
$F \to X$ and that $X_i \times_X F$ is isomorphic to $V_i$.
The functor from left to right is given
by this construction $(V_i, \varphi_{ij}) \mapsto F$.

\medskip\noindent
We omit the verification that these constructions
are mutually quasi-inverse functors. The final statements
(1) and (2) follow from the constructions.
\end{proof}




\begin{multicols}{2}[\section{Autres chapitres}]
\noindent
Préliminaires
\begin{enumerate}
\item \hyperref[introduction-section-phantom]{Introduction}
\item \hyperref[conventions-section-phantom]{Conventions}
\item \hyperref[sets-section-phantom]{Théorie des ensembles}
\item \hyperref[categories-section-phantom]{Catégories}
\item \hyperref[topology-section-phantom]{Topologie}
\item \hyperref[sheaves-section-phantom]{Faisceaux sur les espaces}
\item \hyperref[sites-section-phantom]{Sites et faisceaux}
\item \hyperref[stacks-section-phantom]{Champs}
\item \hyperref[fields-section-phantom]{Corps}
\item \hyperref[algebra-section-phantom]{Algèbre commutative}
\item \hyperref[brauer-section-phantom]{Groupes de Brauer}
\item \hyperref[homology-section-phantom]{Algèbre homologique}
\item \hyperref[derived-section-phantom]{Catégories dérivées}
\item \hyperref[simplicial-section-phantom]{Méthodes simpliciales}
\item \hyperref[more-algebra-section-phantom]{Compléments d'algèbre}
\item \hyperref[smoothing-section-phantom]{Lissage des morphismes d'anneaux}
\item \hyperref[modules-section-phantom]{Faisceaux de modules}
\item \hyperref[sites-modules-section-phantom]{Modules sur les sites}
\item \hyperref[injectives-section-phantom]{Objets injectifs}
\item \hyperref[cohomology-section-phantom]{Cohomologie des faisceaux}
\item \hyperref[sites-cohomology-section-phantom]{Cohomologie sur les sites}
\item \hyperref[dga-section-phantom]{Algèbre différentielle graduée}
\item \hyperref[dpa-section-phantom]{Algèbre à puissances divisées}
\item \hyperref[sdga-section-phantom]{Faisceaux différentiels gradués}
\item \hyperref[hypercovering-section-phantom]{Hyperrecouvrements}
\end{enumerate}
Schémas
\begin{enumerate}
\setcounter{enumi}{25}
\item \hyperref[schemes-section-phantom]{Schémas}
\item \hyperref[constructions-section-phantom]{Constructions de schémas}
\item \hyperref[properties-section-phantom]{Propriétés des schémas}
\item \hyperref[morphisms-section-phantom]{Morphismes de schémas}
\item \hyperref[coherent-section-phantom]{Cohomologie des schémas}
\item \hyperref[divisors-section-phantom]{Diviseurs}
\item \hyperref[limits-section-phantom]{Limites de schémas}
\item \hyperref[varieties-section-phantom]{Variétés}
\item \hyperref[topologies-section-phantom]{Topologies sur les schémas}
\item \hyperref[descent-section-phantom]{Descente}
\item \hyperref[perfect-section-phantom]{Catégories dérivées de schémas}
\item \hyperref[more-morphisms-section-phantom]{Compléments sur les morphismes}
\item \hyperref[flat-section-phantom]{Compléments sur la platitude}
\item \hyperref[groupoids-section-phantom]{Schémas en groupoïdes}
\item \hyperref[more-groupoids-section-phantom]{Compléments sur les schémas en groupoïdes}
\item \hyperref[etale-section-phantom]{Morphismes étales de schémas}
\end{enumerate}
Thèmes de théorie des schémas
\begin{enumerate}
\setcounter{enumi}{41}
\item \hyperref[chow-section-phantom]{Homologie de Chow}
\item \hyperref[intersection-section-phantom]{Théorie de l'intersection}
\item \hyperref[pic-section-phantom]{Schémas de Picard des courbes}
\item \hyperref[weil-section-phantom]{Théories cohomologiques de Weil}
\item \hyperref[adequate-section-phantom]{Modules adéquats}
\item \hyperref[dualizing-section-phantom]{Complexes dualisants}
\item \hyperref[duality-section-phantom]{Dualité pour les schémas}
\item \hyperref[discriminant-section-phantom]{Discriminants et différentes}
\item \hyperref[derham-section-phantom]{Cohomologie de de Rham}
\item \hyperref[local-cohomology-section-phantom]{Cohomologie locale}
\item \hyperref[algebraization-section-phantom]{Géométrie algébrique et formelle}
\item \hyperref[curves-section-phantom]{Courbes algébriques}
\item \hyperref[resolve-section-phantom]{Résolution des surfaces}
\item \hyperref[models-section-phantom]{Réduction semi-stable}
\item \hyperref[functors-section-phantom]{Foncteurs et morphismes}
\item \hyperref[equiv-section-phantom]{Catégories dérivées de variétés}
\item \hyperref[pione-section-phantom]{Groupes fondamentaux de schémas}
\item \hyperref[etale-cohomology-section-phantom]{Cohomologie étale}
\item \hyperref[crystalline-section-phantom]{Cohomologie cristalline}
\item \hyperref[proetale-section-phantom]{Cohomologie pro-étale}
\item \hyperref[relative-cycles-section-phantom]{Cycles relatifs}
\item \hyperref[more-etale-section-phantom]{Compléments de cohomologie étale}
\item \hyperref[trace-section-phantom]{La formule des traces}
\end{enumerate}
Espaces algébriques
\begin{enumerate}
\setcounter{enumi}{64}
\item \hyperref[spaces-section-phantom]{Espaces algébriques}
\item \hyperref[spaces-properties-section-phantom]{Propriétés des espaces algébriques}
\item \hyperref[spaces-morphisms-section-phantom]{Morphismes d'espaces algébriques}
\item \hyperref[decent-spaces-section-phantom]{Espaces algébriques décents}
\item \hyperref[spaces-cohomology-section-phantom]{Cohomologie des espaces algébriques}
\item \hyperref[spaces-limits-section-phantom]{Limites d'espaces algébriques}
\item \hyperref[spaces-divisors-section-phantom]{Diviseurs sur les espaces algébriques}
\item \hyperref[spaces-over-fields-section-phantom]{Espaces algébriques sur des corps}
\item \hyperref[spaces-topologies-section-phantom]{Topologies sur les espaces algébriques}
\item \hyperref[spaces-descent-section-phantom]{Descente et espaces algébriques}
\item \hyperref[spaces-perfect-section-phantom]{Catégories dérivées d'espaces}
\item \hyperref[spaces-more-morphisms-section-phantom]{Compléments sur les morphismes d'espaces}
\item \hyperref[spaces-flat-section-phantom]{Platitude sur les espaces algébriques}
\item \hyperref[spaces-groupoids-section-phantom]{Groupoïdes dans les espaces algébriques}
\item \hyperref[spaces-more-groupoids-section-phantom]{Compléments sur les groupoïdes dans les espaces}
\item \hyperref[bootstrap-section-phantom]{Amorçage}
\item \hyperref[spaces-pushouts-section-phantom]{Sommes amalgamées d'espaces algébriques}
\end{enumerate}
Thèmes de géométrie
\begin{enumerate}
\setcounter{enumi}{81}
\item \hyperref[spaces-chow-section-phantom]{Groupes de Chow des espaces}
\item \hyperref[groupoids-quotients-section-phantom]{Quotients de groupoïdes}
\item \hyperref[spaces-more-cohomology-section-phantom]{Compléments sur la cohomologie des espaces}
\item \hyperref[spaces-simplicial-section-phantom]{Espaces simpliciaux}
\item \hyperref[spaces-duality-section-phantom]{Dualité pour les espaces}
\item \hyperref[formal-spaces-section-phantom]{Espaces algébriques formels}
\item \hyperref[restricted-section-phantom]{Algébrisation des espaces formels}
\item \hyperref[spaces-resolve-section-phantom]{Résolution des surfaces revisitée}
\end{enumerate}
Théorie des déformations
\begin{enumerate}
\setcounter{enumi}{89}
\item \hyperref[formal-defos-section-phantom]{Théorie formelle des déformations}
\item \hyperref[defos-section-phantom]{Théorie des déformations}
\item \hyperref[cotangent-section-phantom]{Le complexe cotangent}
\item \hyperref[examples-defos-section-phantom]{Problèmes de déformation}
\end{enumerate}
Champs algébriques
\begin{enumerate}
\setcounter{enumi}{93}
\item \hyperref[algebraic-section-phantom]{Champs algébriques}
\item \hyperref[examples-stacks-section-phantom]{Exemples de champs}
\item \hyperref[stacks-sheaves-section-phantom]{Faisceaux sur les champs algébriques}
\item \hyperref[criteria-section-phantom]{Critères de représentabilité}
\item \hyperref[artin-section-phantom]{Axiomes d'Artin}
\item \hyperref[quot-section-phantom]{Espaces de Quot et de Hilbert}
\item \hyperref[stacks-properties-section-phantom]{Propriétés des champs algébriques}
\item \hyperref[stacks-morphisms-section-phantom]{Morphismes de champs algébriques}
\item \hyperref[stacks-limits-section-phantom]{Limites de champs algébriques}
\item \hyperref[stacks-cohomology-section-phantom]{Cohomologie des champs algébriques}
\item \hyperref[stacks-perfect-section-phantom]{Catégories dérivées de champs}
\item \hyperref[stacks-introduction-section-phantom]{Introduction aux champs algébriques}
\item \hyperref[stacks-more-morphisms-section-phantom]{Compléments sur les morphismes de champs}
\item \hyperref[stacks-geometry-section-phantom]{La géométrie des champs}
\end{enumerate}
Thèmes de théorie des espaces de modules
\begin{enumerate}
\setcounter{enumi}{107}
\item \hyperref[moduli-section-phantom]{Champs de modules}
\item \hyperref[moduli-curves-section-phantom]{Espaces de modules de courbes}
\end{enumerate}
Divers
\begin{enumerate}
\setcounter{enumi}{109}
\item \hyperref[examples-section-phantom]{Exemples}
\item \hyperref[exercises-section-phantom]{Exercices}
\item \hyperref[guide-section-phantom]{Guide bibliographique}
\item \hyperref[desirables-section-phantom]{Desiderata}
\item \hyperref[coding-section-phantom]{Style de programmation}
\item \hyperref[obsolete-section-phantom]{Obsolète}
\item \hyperref[fdl-section-phantom]{Licence de documentation libre GNU}
\item \hyperref[index-section-phantom]{Index généré automatiquement}
\end{enumerate}
\end{multicols}


\bibliography{my}
\bibliographystyle{amsalpha}


\end{document}
